\documentclass[12pt]{article}

% --- pdfLaTeX совместимость (кириллица) ---
\usepackage[T2A]{fontenc}
\usepackage[utf8]{inputenc}
\DeclareUnicodeCharacter{03C7}{\ensuremath{\chi}}
\usepackage{placeins}

\usepackage[russian,ukrainian]{babel}
\usepackage{fancyvrb}
\usepackage{microtype}
\usepackage{geometry}
\geometry{margin=25mm}


\usepackage{tikz}
\usetikzlibrary{arrows.meta,positioning,shapes.geometric}
\usepackage{xcolor}
\usepackage{graphicx} % для \resizebox

\usepackage{amsmath, amssymb, mathtools}
\usepackage{physics}
\usepackage{hyperref}
\hypersetup{colorlinks=true, linkcolor=blue, urlcolor=blue, citecolor=blue}

\usepackage{enumitem}
\setlist[itemize]{noitemsep, topsep=3pt}
\setlist[enumerate]{noitemsep, topsep=3pt}

% --- Термины проекта (макросы): кириллица НЕ через \mathrm ---
% --- Термины проекта (макросы) — безопасно для pdfLaTeX/T2A ---
% --- Symbol macros (pdfLaTeX-safe)
\newcommand{\Sig}{\ensuremath{\Sigma}}
\newcommand{\chiP}{\ensuremath{\chi}}
\newcommand{\epar}{\ensuremath{e_{\parallel}}}
\newcommand{\Vpar}{\ensuremath{V_{\parallel}}}
\newcommand{\Pitop}{\ensuremath{\Pi_{\top}}}
\newcommand{\vrms}{\ensuremath{v_{\mathrm{rms}}}}
\newcommand{\nbar}{\ensuremath{\bar n}}
\newcommand{\epsG}{\ensuremath{\varepsilon_G}}
\newcommand{\epsA}{\ensuremath{\varepsilon_A}}
\newcommand{\epsV}{\ensuremath{\varepsilon_V}}
\newcommand{\epsVp}{\ensuremath{\varepsilon_{V\parallel}}}
\newcommand{\Jintp}{\ensuremath{J_{\mathrm{int},\parallel}}}
\newcommand{\SE}{\mbox{\textsc{СЭ}}}
\newcommand{\EVO}{\mbox{\textsc{ЭВО}}}
\newcommand{\GP}{\mbox{\textsc{ГП}}}
\newcommand{\KP}{\mbox{\textsc{КП}}}
\newcommand{\VR}{\mbox{\textsc{ВР}}}
\newcommand{\NE}{\mbox{\textsc{НЭ}}}
\newcommand{\NY}{\mbox{\textsc{НЯ}}}
\newcommand{\A}{\mbox{\textsc{А}}} % тут латиница, можно оставить

\usepackage{amsthm}
\theoremstyle{definition}
\newtheorem{definition}{Определение}[section]
\newtheorem{axiom}{Аксиома}[section]
\newtheorem{postulate}{Постулат}[section]
\newtheorem{ruleofthumb}{Рабочее правило}[section]
\newtheorem{hypothesis}{Гипотеза}[section]
\theoremstyle{plain}
\newtheorem{lemma}{Лемма}[section]
\newtheorem{proposition}{Утверждение}[section]
\theoremstyle{remark}
\newtheorem{remark}{Замечание}[section]

\newcommand{\mean}[1]{\left\langle #1 \right\rangle}

\title{Эфирокинетическая теория: рабочий черновик}
\author{Проект «Эфир»}
\date{\today}

\begin{document}
\maketitle

\begin{abstract}
Рабочая версия теории проекта «Эфир». Документ собирает: (i) канон (аксиоматика 1--15),
(ii) математическое ядро (кинетика, граничные условия), (iii) проверяемые следствия и тесты,
(iv) журнал изменений.
\end{abstract}

\tableofcontents
\bigskip

% =========================================================
\section{КАНОН (1--17)}
\subsection*{Статус}
\textbf{Канон:} утверждается руководителем \quad
\textbf{Версия:} v1.5 \quad
\textbf{Дата (Europe/Kyiv):} \today

% ----------------- 1
\subsection{1. Рамки и стиль}
\begin{definition}[Рамка модели]
Вселенная рассматривается как кинетика единой среды. \SE\ (Свободный эфир) — субъект наблюдаемого мира;
\EVO\ — устойчивые вихревые формы организации \SE\ (материя).
Термины ``поле'', ``сила'', ``давление'', ``линии'' допустимы только как макро-ярлыки статистики соударений
и не должны подменять кинетический смысл.
\end{definition}

% ----------------- 2
\subsection{2. Фундаментальные сущности}
\begin{definition}[Амеры]
\A\ — одинаковые сферические, упругие, гладкие частицы. В балансах решают скорость, геометрия и статистика столкновений,
а не ``индивидуальная масса Амера''.
\end{definition}

% ----------------- 3
\subsection{3. Фазовые состояния}
\begin{definition}[Фазы]
Фазы:
\begin{itemize}
\item \textbf{Хаотическая:} \SE\ вне \EVO; \NE\ — хаос внутри \EVO\ (ядро \NY\ и межканальные прослойки).
\item \textbf{Направленная:} только циркуляционные каналы \EVO\ (канальный поток не отождествляется с \NE).
\end{itemize}
\textbf{Запрет:} в \SE\ (как хаотической фазе) нет устойчивых ``потоков''; допустимы лишь распространяющиеся \VR\
как флуктуации импульса и/или плотности. \VR\ не является потоком \SE\ и не является ``субстанцией поля''.
\end{definition}

% ----------------- 4
\subsection{4. Инвариант: единственный физический акт}
\begin{axiom}[Инвариант]
Единственный физический акт — парные упругие соударения Амеров; все процессы (структурирование, перенос возмущений,
взаимодействия, движение и вращение \EVO) реализуются только через них.
Запрещены механизмы ``поглощение/переизлучение Амеров'', ``протекание Амеров через границы'',
``источники поля как отдельная субстанция''.
\end{axiom}

% ----------------- 5
\subsection{5. Целостность: запрет на обгон (хаотическая фаза)}
\begin{postulate}[Условие целостности хаотической фазы]
Запрет на обгон действует только в хаотической фазе (\SE\ и \NE) и является условием устойчивости \EVO:
\[
l \le d_A, \qquad \mean{l} = 0.5\,d_A.
\]
При нарушении этого режима вихревая структура теряет целостность и распадается.
\end{postulate}

\begin{remark}[Каналы]
В каналах (направленная фаза) концентрация ниже, скорости выше; обгоны/отставания возможны внутри потока,
но отсутствует движение против общего направления.
\end{remark}

% ----------------- 6
\subsection{6. Непроницаемость граничных поверхностей и КУБК}
\begin{axiom}[Непроницаемость граничных поверхностей]
Граничные поверхности (\GP) непроницаемы:
\begin{itemize}
\item для Амеров канала — вследствие условия целостности (структура не допускает ``вылет'');
\item для Амеров \SE/\NE\ — вследствие частоты и направленности соударений с Амерами канала на \GP.
\end{itemize}
\end{axiom}

\begin{postulate}[КУБК: кинетическое условие баланса контактов \GP]
В квазистационарном режиме на каждом участке \GP отсутствует устойчивый перекос по \emph{числу} контактов
между амерами хаотической фазы (\SE/\NE) и амерами канала у \GP.

Эквивалентно, после осреднения по времени выполнено согласование \emph{интенсивностей контактов (касания \GP)}:
\[
\Phi^-_{\mathrm{хаос}}(x)=\Phi^+_{\mathrm{канал}}(x),\qquad x\in\Sigma,
\]
где \(\Phi\) — не ``поток \SE\ в объёме'', а локальная величина \emph{частоты/интенсивности касаний поверхности}:
\[
\Phi^-_{\mathrm{хаос}}(x)=\int_{v\cdot n<0}\abs{v\cdot n}\;f_{\mathrm{хаос}}(x,v)\,d^3v,\qquad
\Phi^+_{\mathrm{канал}}(x)=\int_{v\cdot n>0}(v\cdot n)\;f_{\mathrm{канал}}(x,v)\,d^3v.
\]
Здесь \(f_{\mathrm{хаос}}(x,v)\) и \(f_{\mathrm{канал}}(x,v)\) — локальные функции распределения амеров по скоростям
в окрестности \(\Sigma\) со своей стороны поверхности; интегрирование ведётся по соответствующим полупространствам
\(v\cdot n<0\) и \(v\cdot n>0\).

Среди всех режимов канала, совместимых с целостностью \EVO, реализуется режим с \emph{минимально возможной}
локальной концентрацией \(n_{\mathrm{канал}}(x)\) у \GP, при котором указанное согласование сохраняется.
\end{postulate}

\begin{remark}[Различение балансов]
КУБК относится к балансу \emph{числа} контактов. Дрейф/вращение \EVO\ задаются анизотропией
\emph{приращений импульса} в контактах на \GP, а не перекосом по числу контактов.
\end{remark}

% ----------------- 7
\subsection{7. Геометрия \EVO: каналы, тонкий тор и квазицилиндрическое приближение}
\begin{definition}[Канальная структура]
\EVO\ строится из одного или нескольких коаксиальных циркуляционных каналов (направленная фаза) и областей \NE\ (хаос)
между ними и в \NY. У каждого канала две \GP\ (внутренняя и внешняя); внешняя \GP\ внешнего канала — граница \EVO\ с \SE.
\end{definition}

\begin{postulate}[Тонкий тор]
Базовые тороидальные \EVO\ рассматриваются в режиме ``тонкого вихря'': \(r\ll R\),
где \(r\) — малый радиус тора, \(R\) — большой. Это фиксирует локальную однородность условий движения Амеров вдоль большой окружности
и допускает квазицилиндрическое приближение для коротких участков траектории.
\end{postulate}

\begin{remark}[Дополнение (усилитель)]
Условие \(r\ll R\) означает малость изменения локальной кривизны на масштабе \(l\) (длины свободного пробега),
что и делает квазицилиндрическое приближение допустимым.
\end{remark}

% ----------------- 8
\subsection{8. Шнековая организация \GP (динамическая)}
\begin{postulate}[Динамика ``шнека'']
Шнекоподобная \GP\ трактуется как динамический режим: равновесие на \GP поддерживается колебательным процессом границы
относительно среднего положения. Колебания \GP локально ортогональны направлению основного потока в канале.
Период/фаза колебаний связаны (или фазно согласованы) с характерным временем циркуляции в структуре канала
(в частности, с временем обращения вокруг малого кольца тора и/или его гармониками).
Возможны режимы шнека: стоячий и бегущий (в том числе с направлением бега по ходу или против хода канального движения),
в зависимости от распределения внутренних/межканальных \VR\ и внешнего состояния \KP.
\end{postulate}

% ----------------- 9
\subsection{9. Возмущения и кинетический потенциал}
\begin{definition}[\VR]
\VR\ — распространяющиеся в \SE/\NE\ флуктуации импульса и плотности, передающиеся цепочками соударений;
\VR\ не является потоком \SE.
\end{definition}

\begin{definition}[\KP]
\KP\ — характеристика состояния хаотической фазы (распределение скоростей и концентрации Амеров);
градиенты \KP\ задают смещение статистики соударений и приращений импульса на \GP.
\end{definition}

% ----------------- 10
\subsection{10. Движение и вращение \EVO: механизм через \GP и симметрия дрейфа}
\begin{axiom}[Механизм через \GP]
Перемещение/вращение \EVO\ — интеграл микросмещений мест соударений на \GP\ (``Амер хаоса \(\leftrightarrow\) Амер канала на \GP'').
Для движения не требуется ``разность давлений'' как первичная причина: достаточно устойчивого смещения статистики
\textbf{приращений импульса} в контактах на части \(\GP\) $\Rightarrow$ дрейф; момент возникает аналогично при несбалансированности микросмещений по поверхности.
\end{axiom}

\begin{postulate}[Симметрия дрейфа]
В изотропном фоне \(\nabla\KP=0\) и при отсутствии внешних \VR, выделяющих направление, идеализированный симметричный тонкий тор
не имеет выделенного направления для устойчивого дрейфа. Следовательно, дрейф требует нарушения симметрии в статистике
приращений импульса на \GP (внешнего или самосогласованного внутреннего).
\end{postulate}

% ----------------- 11
\subsection{11. Масса}
\begin{definition}[Масса как отклик]
Масса \EVO\ — мера отклика структуры на внешние \VR/\(\nabla\KP\), а не сумма масс ``заключённых Амеров''.
\end{definition}

\begin{ruleofthumb}[Рабочее правило массы]
Масса \(\sim\) суммарной воспринимающей площади \GP\ всех каналов (многоканальность увеличивает ``тяжесть'' при близкой внешней геометрии).
\end{ruleofthumb}

% ----------------- 12
\subsection{12. Заряд и электрическое взаимодействие}
\begin{axiom}[Заряд и электрическое]
Заряд задаётся направленностью/зарученностью (хиральностью) циркуляций в каналах.

\textbf{Нотация (канон):} \(\chi\in\{-1,+1\}\) — дискретный знак хиральности (зарученности) для канала и/или \EVO\ как целого
(по правилу ``паспорта вихря''). Локальные ``поля хиральности'' на \(\Sigma\) не вводятся как базовая сущность; при необходимости
используются как производные рабочие метки.

Взаимодействие: \VR\ двух \EVO\ интерферируют $\Rightarrow$ формируют устойчивые градиенты \KP, которые
(а) ориентируют \EVO\ в конфигурацию притяжения/отталкивания,
(б) через механизм \GP\ реализуют сближение/разлёт.
\end{axiom}

% ----------------- 13
\subsection{13. Фотон, аннигиляция и двуторий как режим ``вращение\,$\to$\,дрейф''}
\begin{axiom}[Фотон и аннигиляция]
Фотон — одноканальный \EVO\ (правый/левый), без заряда; ``световой режим'' связан с согласованной шнековой организацией \GP.
Аннигиляция — топологическая трансформация многоканальной структуры в моноканальную (фотонную) без исчезновения Амеров.
\end{axiom}

\begin{postulate}[Двуторий: самосогласованное нарушение симметрии]
Двуторий (тороидально-спиральная форма \EVO) рассматривается как разрешённый режим, в котором внутренняя топология и шнековая организация \GP
допускают \textbf{самосогласованное} появление устойчивой анизотропии статистики контактов/импульса на разных участках \(\Sigma\).

Механизм трактуется кинетически: собственные \VR\ структуры и их интерференция формируют устойчивую анизотропию состояния хаотической фазы
(градиенты \KP) у разных участков \GP; через интеграл микросмещений мест соударений на \GP это даёт ненулевой дрейф и/или момент
без введения ``полей'' как первичных сущностей.
\end{postulate}

% ----------------- 14
\subsection{14. Гравитация и спин/магнетизм}
\begin{axiom}[Гравитация]
Гравитация — эффект градиента интенсивности/плотности \VR\ в \SE, создаваемого множеством удалённых \EVO;
зависимость \(\sim 1/r^2\) связывается со сферическим распространением \VR\ и интегральной статистикой соударений на \GP.
\end{axiom}

\begin{axiom}[Спин/магнетизм без отдельного поля]
Спин и магнитный момент — проявления топологии и шнековой \GP; внешние \VR/\(\nabla\KP\) дают одновременно продольный дрейф
и вращательный отклик. Магнетизм не вводится как субстанция: ``силовые линии'' — линии хиральной анизотропии статистики соударений в \SE\
при доменной сонаправленности \EVO. Нулевой суммарный заряд не запрещает магнетизм при сохранении осевой хиральности.
\end{axiom}

% ----------------- 15
\subsection{15. Квантование, паспорт вихря и запреты}
\begin{axiom}[Онтологическое квантование]
Квантование не постулируется: устойчивость требует целочисленного замыкания циркуляции/геометрии канала
(самосогласованности статистики соударений на \GP). Простейшая запись для тонкого тора: \(R=n\cdot r\), \(n\in\mathbb{Z}\),
где \(n\) — целое число самосогласованного замыкания (условие устойчивости).
\end{axiom}

\begin{remark}[Трёхканальный вихрь и дробные значения]
Трёхканальный вихрь и знаки зарученности (± по каналам) при нормировке на суммарную воспринимающую площадь дают значения
\(0, \pm\frac{1}{3}, \pm\frac{2}{3}, \pm 1\). Античастица — смена знаков зарученности.
\end{remark}

\begin{definition}[Паспорт вихря (нотация)]
Стрелка с хвостом слева: направление потока в канале вдоль большой окружности.
Ориентация: ``смотреть с конца стрелки'' (по/против часовой).

\textbf{Хиральность (зарученность):} \(\chi\in\{-1,+1\}\).
Паспорт записывается как список знаков \([\chi_1,\chi_2,\dots]\) изнутри наружу и (при необходимости)
общий \(\chi\) объекта как принятую функцию от списка.

Шнек-режим указывать отдельно (тип: стоячий/бегущий; направление бега; фазная связь), без трактовки как ``потока \SE''.
\end{definition}

\begin{axiom}[Жёсткие запреты]
Не вводить: устойчивые потоки \SE, ``давление \SE'' как первичный агент, проникновение Амеров через \GP,
поглощение/переизлучение Амеров, ``виртуальные переносчики'' как физические объекты, ``поле'' как отдельную субстанцию.
Генезис вихрей из \SE\ в рамках канона не обсуждать (\EVO\ считаются существующими).
\end{axiom}

\begin{remark}
Термины ``поле/сила/давление'' допускаются только как макро-ярлыки статистики соударений, без онтологической подмены.
\end{remark}

% ----------------- 16
\subsection{16. Прокси (K, G, A\_ij) как измерительные ярлыки (\KP/\VR), не новые сущности}
\begin{axiom}[Прокси-стандарт тестирования]
В целях сопоставимости расчётов, симуляций и эксперимента допускается использование прокси-наборов
\((K,\,G,\,A_{ij})\) как \emph{измерительных ярлыков} для состояния хаотической фазы в окрестности \(\Sigma\)
(то есть для канонических понятий \(\KP\) и \(\VR\)).

\textbf{Нормативно:}
\begin{enumerate}
\item Прокси \((K,\,G,\,A_{ij})\) \textbf{не являются новыми сущностями}, ``полями'' или ``субстанциями'' и не вводят в модель
никакой дополнительной онтологии.
\item Прокси определяются как функционалы от распределения \(f(x,v)\) амеров хаотической фазы (СЭ/НЭ) у \(\Sigma\),
и служат только для записи измеримых/вычислимых характеристик (скалярной, векторной и тензорной).
\item Любая интерпретация прокси как самостоятельных физических объектов запрещена; прокси — это язык наблюдения,
а не источник причинности. Причинность остаётся кинетической: соударения, КУБК (баланс числа контактов) и анизотропия
приращений импульса на \(\Sigma\).
\end{enumerate}
\end{axiom}

\begin{remark}[Рекомендуемая (не обязательная) нормировка]
Для единообразия вычислений допускается использовать:
\[
K=n\langle v^2\rangle,\qquad G=n\langle v\rangle,\qquad
A_{ij}=P_{ij}-\frac{1}{3}\tr(P)\,\delta_{ij},\qquad P_{ij}=n\langle v_i v_j\rangle.
\]
\end{remark}

% ----------------- 17
\subsection{17. Норматив тестов шнека: \(\chi\)-нечётность продольного отклика (T5--T6)}
\begin{postulate}[\(\chi\)-ручка и отличимость режимов шнека]
Для тестов шнека фиксируется ось объекта \(e_{\parallel}\): \textbf{ось симметрии} тора/двутория, вдоль которой измеряется поступательный дрейф.

\textbf{Нормативно (\(\chi\)-ручка):} при прочих равных (геометрия \(\Sigma\) фиксирована; КУБК соблюдён; режим шнека фиксирован)
\textbf{\(\chi\)-нечётной обязана быть продольная скорость дрейфа}
\[
V_{\parallel}(\chi)\equiv V\cdot e_{\parallel},\qquad V_{\parallel}(\chi)=-V_{\parallel}(-\chi).
\]
То есть смена паспорта хиральности \(\chi\to-\chi\) должна менять знак именно \(V_{\parallel}\).

\textbf{Нормативно (реверс бега шнека):} при фиксированном \(\chi\) и прочих равных смена направления бега шнека
обязана менять знак \(V_{\parallel}\):
\[
V_{\parallel}(\text{бег}\,\rightarrow)=-V_{\parallel}(\text{бег}\,\leftarrow).
\]
Если режим шнека стоячий, то \(V_{\parallel}=0\) является нормативной проверкой (при отсутствии внешнего нарушителя).
\end{postulate}

\begin{remark}[Следствие для моделирования]
В любой замыкательной схеме (в т.ч. через прокси \(G, A_{ij}\) на \(\Sigma\)) \(\chi\) должна входить так, чтобы выделяемая
\(\chi\)-нечётная часть отклика давала вклад в \(V_{\parallel}\), а не подменялась ``скрытым потоком \SE'' или
асимметрией числа контактов (что запрещено логикой КУБК).
\end{remark}

% =========================================================
% sec_axioms_A.tex — A/Аксиоматика (вставка) — АКТУАЛЬНАЯ РЕДАКЦИЯ
% Назначение: строгие определения/аксиомы/постулаты, на которые
% ссылаются разделы математики (B) и верификации (C).
% Подключение: % =========================================================
% sec_axioms_A.tex — A/Аксиоматика (вставка) — АКТУАЛЬНАЯ РЕДАКЦИЯ
% Назначение: строгие определения/аксиомы/постулаты, на которые
% ссылаются разделы математики (B) и верификации (C).
% Подключение: % =========================================================
% sec_axioms_A.tex — A/Аксиоматика (вставка) — АКТУАЛЬНАЯ РЕДАКЦИЯ
% Назначение: строгие определения/аксиомы/постулаты, на которые
% ссылаются разделы математики (B) и верификации (C).
% Подключение: \input{sec_axioms_A}
% Совместимость: article, pdfLaTeX, T2A, babel (ru/uk)
% =========================================================

\section{Аксиоматика и определения A (расширение канона)}

\subsection*{Статус}
\textbf{Ответственный:} Сотрудник A (аксиоматика) \quad
\textbf{Статус:} рабочая вставка \quad
\textbf{Версия:} A-0.2 \quad
\textbf{Дата (Europe/Kyiv):} \today

% ---------------------------------------------------------
\subsection{A.1. Термины и фазовые состояния}

\begin{definition}[Фазовые состояния эфира]
В модели различаются два фазовых состояния движения Амеров:
\begin{itemize}
\item \textbf{Хаотическая фаза:} статистически хаотическое движение Амеров.
Вне \EVO\ хаотическая фаза называется \SE\ (Свободный эфир), а внутри \EVO\ — \NE\ (Несвободный эфир),
включая центральное ядро \NY\ и межканальные прослойки.
\item \textbf{Направленная фаза:} упорядоченное циркуляционное движение, существующее \emph{только} в каналах \EVO.
Направленное движение в каналах \textbf{не} называется \NE.
\end{itemize}
\end{definition}

\begin{axiom}[Запрет устойчивых потоков \SE]
В \SE\ не допускаются устойчивые потоки как первичная онтология.
Допускаются только распространяющиеся \VR\ как флуктуации импульса/плотности,
передаваемые цепочками парных соударений.
\end{axiom}

% ---------------------------------------------------------
\subsection{A.2. Геометрия \EVO: каналы и граничные поверхности}

\begin{definition}[Каналы и граничные поверхности (\GP)]
\EVO\ состоит из одного или нескольких коаксиальных циркуляционных каналов (направленная фаза)
и областей \NE\ (хаос) между ними и в \NY.
Каждый канал имеет две \GP: внутреннюю и внешнюю.
Внешняя \GP внешнего канала является внешней \GP всего \EVO\ (границей \EVO\ с \SE).
\end{definition}

\begin{remark}[Шнековая организация \GP]
Шнекоподобная форма \GP понимается как согласованная пространственно-временная организация контактов,
которая может быть бегущей или стоячей (см. A.6) и влияет на статистику соударений на \GP.
\end{remark}

% ---------------------------------------------------------
\subsection{A.3. Инвариант и жёсткие запреты}

\begin{axiom}[Единственный физический акт]
Единственный физический акт модели — парные упругие соударения Амеров.
Все процессы (структурирование, перенос \VR, взаимодействия, движение/вращение \EVO) реализуются только через них.
\end{axiom}

\begin{axiom}[Жёсткие запреты модели]
Запрещается вводить как первичные сущности:
\begin{itemize}
\item устойчивые потоки \SE;
\item ``давление \SE'' как первичный агент;
\item проникновение Амеров через \GP;
\item поглощение/переизлучение Амеров;
\item ``виртуальные переносчики'' как физические объекты;
\item ``поле'' как отдельную субстанцию (термин допустим только как макро-ярлык).
\end{itemize}
Генезис вихрей из \SE\ в рамках канона не рассматривается (\EVO\ считаются существующими).
\end{axiom}

% ---------------------------------------------------------
\subsection{A.4. Условие целостности хаотической фазы (запрет на обгон)}

\begin{postulate}[Целостность хаотической фазы у \GP]
В хаотической фазе (\SE\ и \NE) у \GP поддерживается режим, обеспечивающий целостность \EVO:
\[
l_{\max}\le d_A, \qquad \mean{l}=0.5\,d_A,
\]
где \(l\) — длина свободного пробега Амера, \(d_A\) — диаметр Амера.
При нарушении режима целостности вихревая структура теряет устойчивость и распадается.
\end{postulate}

\begin{remark}[Каналы]
В каналах (направленная фаза) концентрация ниже, скорости выше;
обгоны/отставания возможны внутри потока, но отсутствует движение против общего направления.
\end{remark}

% ---------------------------------------------------------
\subsection{A.5. КУБК: контактный баланс и минимальная концентрация канала}

\begin{postulate}[КУБК: баланс \emph{числа} контактов на \GP]
В квазистационарном режиме на каждом участке \GP отсутствует устойчивый перекос по \emph{числу} контактов
между амерами хаотической фазы (\SE/\NE) и амерами канала у \GP.
Эквивалентно, после осреднения по времени выполняется согласование \emph{интенсивностей контактов (касания \GP)}:
\[
\Phi^-_{\mathrm{хаос}}(x)=\Phi^+_{\mathrm{канал}}(x),\qquad x\in\Sigma,
\]
где \(\Sigma\) — соответствующая \GP, \(n\) — внешняя единичная нормаль к \(\Sigma\),
а \(\Phi\) не является ``потоком \SE\ в объёме'', а есть локальная величина \emph{частоты/интенсивности касаний поверхности}:
\[
\Phi^-_{\mathrm{хаос}}(x)=\int_{v\cdot n<0}\abs{v\cdot n}\;f_{\mathrm{хаос}}(x,v)\,d^3v,\qquad
\Phi^+_{\mathrm{канал}}(x)=\int_{v\cdot n>0}(v\cdot n)\;f_{\mathrm{канал}}(x,v)\,d^3v.
\]
Здесь \(f_{\mathrm{хаос}}(x,v)\) и \(f_{\mathrm{канал}}(x,v)\) — локальные функции распределения амеров по скоростям
в окрестности \(\Sigma\) со своей стороны поверхности; интегрирование ведётся по соответствующим полупространствам
\(v\cdot n<0\) и \(v\cdot n>0\).
\end{postulate}

\begin{postulate}[Принцип минимальной концентрации канала у \GP]
Среди всех режимов канала, совместимых с целостностью \EVO, реализуется режим с \emph{минимально возможной}
локальной концентрацией \(n_{\mathrm{канал}}(x)\) у \GP, при котором сохраняется КУБК.
\end{postulate}

\begin{remark}[Различение балансов]
КУБК фиксирует баланс \emph{числа} контактов.
Дрейф/вращение определяются не этим балансом, а анизотропией \emph{приращений импульса} в контактах на \GP.
\end{remark}

% ---------------------------------------------------------
\subsection{A.6. Динамическая природа ``шнека''}

\begin{postulate}[Шнек как колебательный режим \GP]
Равновесие на \GP поддерживается колебательным процессом границы относительно среднего положения.
Колебания \GP локально ортогональны направлению основного потока в канале.
Период/фаза колебаний связаны (или фазно согласованы) с характерным временем циркуляции в структуре канала,
в частности с временем обращения вокруг малого кольца тора и/или его гармониками.
Возможны режимы: стоячий шнек и бегущий шнек (по ходу или против хода канального движения),
в зависимости от распределения внутренних/межканальных \VR\ и состояния \KP.
\end{postulate}

% ---------------------------------------------------------
\subsection{A.7. Симметрия дрейфа, хиральность и двуторий}

\begin{postulate}[Симметрия дрейфа в изотропном фоне]
В изотропном фоне (\(\nabla\KP=0\)) и при отсутствии внешних \VR, выделяющих направление,
идеализированный симметричный ``тонкий'' тор не имеет выделенного направления для устойчивого дрейфа.
Следовательно, дрейф требует нарушения симметрии статистики \emph{приращений импульса} на \GP
(внешнего или самосогласованного внутреннего).
\end{postulate}

\begin{definition}[Хиральность (зарученность)]
Хиральность \(\chi\) — дискретный паспортный знак зарученности циркуляции канала/объекта:
\[
\chi\in\{-1,+1\}.
\]
\(\chi\) используется как метка конфигурации (паспорт вихря) и не вводится как отдельная физическая субстанция.
\end{definition}

\begin{definition}[Двуторий]
Двуторий — тороидально-спиральная форма \EVO, возникающая как самосогласованная трансформация ``тонкого'' вихря,
когда собственные \VR\ и связанная с ними структура \KP\ внутри \EVO\ приводят к устойчивому нарушению симметрии
статистики контактов/приращений импульса на \GP.
\end{definition}

\begin{postulate}[Двуторий: механизм ``вращение \(\to\) дрейф'']
В двутории внутренняя топология и шнековая организация \GP обеспечивают преобразование вращательных мод
в поступательный дрейф вдоль оси симметрии.
Механизм трактуется кинетически: интерференция собственных \VR\ формирует устойчивую анизотропию состояния хаотической фазы
(градиенты \KP) у разных участков \GP; интеграл микросмещений мест соударений на \GP даёт ненулевой дрейф и/или момент.
\end{postulate}

\begin{remark}
Формализация дрейфа/момента через поверхностные интегралы от анизотропии импульсного обмена \(j(x,t)\)
является предметом раздела B/Математика.
\end{remark}

% ---------------------------------------------------------
\subsection{A.8. Прокси (K, G, A\_ij): измерительные ярлыки (норматив)}

\begin{axiom}[Прокси как язык измерений, не сущности]
В целях сопоставимости расчётов, симуляций и эксперимента допускается использование прокси-наборов
\((K,\,G,\,A_{ij})\) как \emph{измерительных ярлыков} для состояния хаотической фазы в окрестности \(\Sigma\)
(т.е. для канонических понятий \(\KP\) и \(\VR\)).

\textbf{Нормативно:} прокси не являются новыми сущностями и не добавляют причинности в модель.
Причинность остаётся кинетической: соударения, КУБК (баланс числа контактов) и анизотропия приращений импульса на \(\Sigma\).
\end{axiom}

\begin{remark}[Запретная лексика для прокси]
Запрещены формулировки причинности вида: ``\(G\) действует/толкает'',
``тензор \(A_{ij}\) создаёт силу'', ``поле \(K\) ускоряет'' и т.п.
Допустимо только: ``прокси \(G\)/\(A_{ij}\) наблюдается/оценён'' и ``используется как измерительная метка''.
\end{remark}

\begin{remark}[Рекомендуемая нормировка (не обязательна)]
Допускается использовать единый вид для первых моментов распределения:
\[
K=n\langle v^2\rangle,\qquad G=n\langle v\rangle,\qquad
A_{ij}=P_{ij}-\frac{1}{3}\tr(P)\,\delta_{ij},\qquad P_{ij}=n\langle v_i v_j\rangle.
\]
\end{remark}

% ---------------------------------------------------------
\subsection{A.9. Нормы тестов шнека: ось e\_\parallel и разведение ``V=0''}

\begin{postulate}[Канонический выбор оси \(e_{\parallel}\)]
Для тестов шнека вводится ось объекта \(e_{\parallel}\) — \textbf{ось симметрии} тора/двутория,
вдоль которой определяется продольный дрейф.

\textbf{Нормативно:} \(e_{\parallel}\) должен быть однозначно воспроизводим из геометрии объекта и его паспорта,
без апелляции к внешним ``полям''.
Для тора/двутория ориентация \(e_{\parallel}\) согласуется с принятой параметризацией \(X(\varphi,\theta)\) поверхности и паспортом вихря:
касательное направление \(\partial_{\varphi}X\) соответствует циркуляции по большой окружности, а \(e_{\parallel}\) — оси симметрии выбранной системы координат.
\end{postulate}

\begin{postulate}[Разведение норм ``\(V=0\)'' и нормы по \(V_{\parallel}\)]
В тестах используются две разные нормы ``нуля'', их нельзя смешивать:
\begin{enumerate}
\item \textbf{Охранная норма (T2/T6):} запись ``\(V=0\)'' означает \textbf{нулевой вектор} скорости дрейфа,
т.е. ноль всех компонент:
\[
V=0 \iff V_x=V_y=V_z=0.
\]
\item \textbf{Норма шнека (T5--T6):} нормативной величиной является \textbf{продольная компонента}
\[
V_{\parallel}\equiv V\cdot e_{\parallel}.
\]
При прочих равных \(\chi\)-нечётной обязана быть именно \(V_{\parallel}\):
\[
V_{\parallel}(\chi)=-V_{\parallel}(-\chi).
\]
При фиксированном \(\chi\) реверс направления бега шнека обязан реверсировать знак \(V_{\parallel}\).
Для стоячего шнека нормативно \(V_{\parallel}=0\) (при отсутствии внешнего нарушителя).
\end{enumerate}
\end{postulate}

\begin{remark}
Смысл нормы: \(\chi\) должна входить в замыкания так, чтобы \(\chi\)-нечётная часть отклика давала вклад в \(V_{\parallel}\),
а не подменялась ``скрытым потоком \SE'' или асимметрией числа контактов.
\end{remark}

\section{Аксиоматическая рамка текущего шага: локальный патч как временное сужение задачи}
\label{sec:a_current_step}

\subsection{Глобальная цель и смысл временного сужения}

Цель текущего шага не состоит в замене глобальной тороидальной или двуториальной
геометрии ЭВО локальным цилиндрическим описанием.
Глобальная цель остаётся прежней:
выявить кинетический механизм дрейфа и ротации ЭВО через контактную статистику
и анизотропию приращений импульса на граничной поверхности $\Sigma$.

Квазицилиндрическое приближение вводится только как временное локальное сужение задачи.
Оно используется там, где режим тонкого тора $r \ll R$ позволяет считать короткий участок
большой окружности локально цилиндрическим.
Тем самым локальный патч $\Sigma_{\mathrm{loc}} \subset \Sigma$
служит не новой онтологией, а увеличительным языком записи,
удобным для выделения локальных механизмов на ГП.

\subsection{Что именно сохраняется неизменным}

В аксиоматическом смысле сохраняются следующие положения:

\begin{enumerate}
\item ЭВО глобально остаётся тороидальной или двуториальной структурой.
\item КУБК относится только к балансу числа контактов на $\Sigma$.
\item Дрейф и вращение допускаются только через анизотропию $\Delta p$ на $\Sigma$.
\item Шнек остаётся динамическим режимом ГП, а не отдельной сущностью и не потоком СЭ.
\item Локальный патч не заменяет глобальный интеграл по $\Sigma$, а только помогает его разложить.
\end{enumerate}

Следовательно, переход к $\Sigma_{\mathrm{loc}}$ не является сменой физической картины,
а является временным усилением разрешающей способности математического языка.

\subsection{Почему непрерывная q-схема оказалась переходной}

На предыдущем этапе непрерывный параметр $q$ использовался как язык локальной намотки.
Однако такой язык оказался недостаточным для сохранения физики шнекового процесса,
поскольку не различал явно:

\begin{itemize}
\item траектории Амеров в канале;
\item шнековую моду граничной поверхности;
\item дискретную связность шнековой организации с канальной навивкой.
\end{itemize}

Поэтому в текущем шаге физический приоритет непрерывного $q$ снимается.
Если параметр $q$ и сохраняется, то только как вспомогательный эффективный
локальный параметр записи $q_{\mathrm{eff}}$.

\subsection{Новая кандидатная рамка текущего шага}

В рабочем контуре проекта принимается кандидатная схема:

\[
m_A \in \mathbb{Z}_{>0},
\qquad
m_S \in \{1,2,4\},
\]

где $m_A$ задаёт квантовое число канальной навивки,
а $m_S$ задаёт базовую пространственную моду шнека на ГП.

Базовыми режимами считаются:
\[
S1,\quad S2,\quad S4,
\]
а все прочие режимы допускаются только как производные от них.

В локальном патч-языке должны различаться два касательных направления:
\[
h_A \quad \text{и} \quad h_S,
\]
где $h_A$ есть директор траектории Амеров,
а $h_S$ --- директор шнековой моды ГП.
Шнек не отождествляется с направлением канальной траектории.

\subsection{Интерпретационная граница}

Текущий шаг проекта следует понимать как последовательность:

\[
\text{глобальная цель}
\;\to\;
\text{тонкий тор}
\;\to\;
\Sigma_{\mathrm{loc}}
\;\to\;
(h_A,h_S)
\;\to\;
(m_A,m_S,S1/S2/S4)
\;\to\;
\text{возврат к глобальному интегралу по } \Sigma.
\]

Тем самым локальный акцент на квазицилиндрическом патче не замещает канон,
а служит промежуточным этапом прояснения механизма.
Именно в таком виде данный шаг должен отражаться в витрине и в журнале изменений.

% Совместимость: article, pdfLaTeX, T2A, babel (ru/uk)
% =========================================================

\section{Аксиоматика и определения A (расширение канона)}

\subsection*{Статус}
\textbf{Ответственный:} Сотрудник A (аксиоматика) \quad
\textbf{Статус:} рабочая вставка \quad
\textbf{Версия:} A-0.2 \quad
\textbf{Дата (Europe/Kyiv):} \today

% ---------------------------------------------------------
\subsection{A.1. Термины и фазовые состояния}

\begin{definition}[Фазовые состояния эфира]
В модели различаются два фазовых состояния движения Амеров:
\begin{itemize}
\item \textbf{Хаотическая фаза:} статистически хаотическое движение Амеров.
Вне \EVO\ хаотическая фаза называется \SE\ (Свободный эфир), а внутри \EVO\ — \NE\ (Несвободный эфир),
включая центральное ядро \NY\ и межканальные прослойки.
\item \textbf{Направленная фаза:} упорядоченное циркуляционное движение, существующее \emph{только} в каналах \EVO.
Направленное движение в каналах \textbf{не} называется \NE.
\end{itemize}
\end{definition}

\begin{axiom}[Запрет устойчивых потоков \SE]
В \SE\ не допускаются устойчивые потоки как первичная онтология.
Допускаются только распространяющиеся \VR\ как флуктуации импульса/плотности,
передаваемые цепочками парных соударений.
\end{axiom}

% ---------------------------------------------------------
\subsection{A.2. Геометрия \EVO: каналы и граничные поверхности}

\begin{definition}[Каналы и граничные поверхности (\GP)]
\EVO\ состоит из одного или нескольких коаксиальных циркуляционных каналов (направленная фаза)
и областей \NE\ (хаос) между ними и в \NY.
Каждый канал имеет две \GP: внутреннюю и внешнюю.
Внешняя \GP внешнего канала является внешней \GP всего \EVO\ (границей \EVO\ с \SE).
\end{definition}

\begin{remark}[Шнековая организация \GP]
Шнекоподобная форма \GP понимается как согласованная пространственно-временная организация контактов,
которая может быть бегущей или стоячей (см. A.6) и влияет на статистику соударений на \GP.
\end{remark}

% ---------------------------------------------------------
\subsection{A.3. Инвариант и жёсткие запреты}

\begin{axiom}[Единственный физический акт]
Единственный физический акт модели — парные упругие соударения Амеров.
Все процессы (структурирование, перенос \VR, взаимодействия, движение/вращение \EVO) реализуются только через них.
\end{axiom}

\begin{axiom}[Жёсткие запреты модели]
Запрещается вводить как первичные сущности:
\begin{itemize}
\item устойчивые потоки \SE;
\item ``давление \SE'' как первичный агент;
\item проникновение Амеров через \GP;
\item поглощение/переизлучение Амеров;
\item ``виртуальные переносчики'' как физические объекты;
\item ``поле'' как отдельную субстанцию (термин допустим только как макро-ярлык).
\end{itemize}
Генезис вихрей из \SE\ в рамках канона не рассматривается (\EVO\ считаются существующими).
\end{axiom}

% ---------------------------------------------------------
\subsection{A.4. Условие целостности хаотической фазы (запрет на обгон)}

\begin{postulate}[Целостность хаотической фазы у \GP]
В хаотической фазе (\SE\ и \NE) у \GP поддерживается режим, обеспечивающий целостность \EVO:
\[
l_{\max}\le d_A, \qquad \mean{l}=0.5\,d_A,
\]
где \(l\) — длина свободного пробега Амера, \(d_A\) — диаметр Амера.
При нарушении режима целостности вихревая структура теряет устойчивость и распадается.
\end{postulate}

\begin{remark}[Каналы]
В каналах (направленная фаза) концентрация ниже, скорости выше;
обгоны/отставания возможны внутри потока, но отсутствует движение против общего направления.
\end{remark}

% ---------------------------------------------------------
\subsection{A.5. КУБК: контактный баланс и минимальная концентрация канала}

\begin{postulate}[КУБК: баланс \emph{числа} контактов на \GP]
В квазистационарном режиме на каждом участке \GP отсутствует устойчивый перекос по \emph{числу} контактов
между амерами хаотической фазы (\SE/\NE) и амерами канала у \GP.
Эквивалентно, после осреднения по времени выполняется согласование \emph{интенсивностей контактов (касания \GP)}:
\[
\Phi^-_{\mathrm{хаос}}(x)=\Phi^+_{\mathrm{канал}}(x),\qquad x\in\Sigma,
\]
где \(\Sigma\) — соответствующая \GP, \(n\) — внешняя единичная нормаль к \(\Sigma\),
а \(\Phi\) не является ``потоком \SE\ в объёме'', а есть локальная величина \emph{частоты/интенсивности касаний поверхности}:
\[
\Phi^-_{\mathrm{хаос}}(x)=\int_{v\cdot n<0}\abs{v\cdot n}\;f_{\mathrm{хаос}}(x,v)\,d^3v,\qquad
\Phi^+_{\mathrm{канал}}(x)=\int_{v\cdot n>0}(v\cdot n)\;f_{\mathrm{канал}}(x,v)\,d^3v.
\]
Здесь \(f_{\mathrm{хаос}}(x,v)\) и \(f_{\mathrm{канал}}(x,v)\) — локальные функции распределения амеров по скоростям
в окрестности \(\Sigma\) со своей стороны поверхности; интегрирование ведётся по соответствующим полупространствам
\(v\cdot n<0\) и \(v\cdot n>0\).
\end{postulate}

\begin{postulate}[Принцип минимальной концентрации канала у \GP]
Среди всех режимов канала, совместимых с целостностью \EVO, реализуется режим с \emph{минимально возможной}
локальной концентрацией \(n_{\mathrm{канал}}(x)\) у \GP, при котором сохраняется КУБК.
\end{postulate}

\begin{remark}[Различение балансов]
КУБК фиксирует баланс \emph{числа} контактов.
Дрейф/вращение определяются не этим балансом, а анизотропией \emph{приращений импульса} в контактах на \GP.
\end{remark}

% ---------------------------------------------------------
\subsection{A.6. Динамическая природа ``шнека''}

\begin{postulate}[Шнек как колебательный режим \GP]
Равновесие на \GP поддерживается колебательным процессом границы относительно среднего положения.
Колебания \GP локально ортогональны направлению основного потока в канале.
Период/фаза колебаний связаны (или фазно согласованы) с характерным временем циркуляции в структуре канала,
в частности с временем обращения вокруг малого кольца тора и/или его гармониками.
Возможны режимы: стоячий шнек и бегущий шнек (по ходу или против хода канального движения),
в зависимости от распределения внутренних/межканальных \VR\ и состояния \KP.
\end{postulate}

% ---------------------------------------------------------
\subsection{A.7. Симметрия дрейфа, хиральность и двуторий}

\begin{postulate}[Симметрия дрейфа в изотропном фоне]
В изотропном фоне (\(\nabla\KP=0\)) и при отсутствии внешних \VR, выделяющих направление,
идеализированный симметричный ``тонкий'' тор не имеет выделенного направления для устойчивого дрейфа.
Следовательно, дрейф требует нарушения симметрии статистики \emph{приращений импульса} на \GP
(внешнего или самосогласованного внутреннего).
\end{postulate}

\begin{definition}[Хиральность (зарученность)]
Хиральность \(\chi\) — дискретный паспортный знак зарученности циркуляции канала/объекта:
\[
\chi\in\{-1,+1\}.
\]
\(\chi\) используется как метка конфигурации (паспорт вихря) и не вводится как отдельная физическая субстанция.
\end{definition}

\begin{definition}[Двуторий]
Двуторий — тороидально-спиральная форма \EVO, возникающая как самосогласованная трансформация ``тонкого'' вихря,
когда собственные \VR\ и связанная с ними структура \KP\ внутри \EVO\ приводят к устойчивому нарушению симметрии
статистики контактов/приращений импульса на \GP.
\end{definition}

\begin{postulate}[Двуторий: механизм ``вращение \(\to\) дрейф'']
В двутории внутренняя топология и шнековая организация \GP обеспечивают преобразование вращательных мод
в поступательный дрейф вдоль оси симметрии.
Механизм трактуется кинетически: интерференция собственных \VR\ формирует устойчивую анизотропию состояния хаотической фазы
(градиенты \KP) у разных участков \GP; интеграл микросмещений мест соударений на \GP даёт ненулевой дрейф и/или момент.
\end{postulate}

\begin{remark}
Формализация дрейфа/момента через поверхностные интегралы от анизотропии импульсного обмена \(j(x,t)\)
является предметом раздела B/Математика.
\end{remark}

% ---------------------------------------------------------
\subsection{A.8. Прокси (K, G, A\_ij): измерительные ярлыки (норматив)}

\begin{axiom}[Прокси как язык измерений, не сущности]
В целях сопоставимости расчётов, симуляций и эксперимента допускается использование прокси-наборов
\((K,\,G,\,A_{ij})\) как \emph{измерительных ярлыков} для состояния хаотической фазы в окрестности \(\Sigma\)
(т.е. для канонических понятий \(\KP\) и \(\VR\)).

\textbf{Нормативно:} прокси не являются новыми сущностями и не добавляют причинности в модель.
Причинность остаётся кинетической: соударения, КУБК (баланс числа контактов) и анизотропия приращений импульса на \(\Sigma\).
\end{axiom}

\begin{remark}[Запретная лексика для прокси]
Запрещены формулировки причинности вида: ``\(G\) действует/толкает'',
``тензор \(A_{ij}\) создаёт силу'', ``поле \(K\) ускоряет'' и т.п.
Допустимо только: ``прокси \(G\)/\(A_{ij}\) наблюдается/оценён'' и ``используется как измерительная метка''.
\end{remark}

\begin{remark}[Рекомендуемая нормировка (не обязательна)]
Допускается использовать единый вид для первых моментов распределения:
\[
K=n\langle v^2\rangle,\qquad G=n\langle v\rangle,\qquad
A_{ij}=P_{ij}-\frac{1}{3}\tr(P)\,\delta_{ij},\qquad P_{ij}=n\langle v_i v_j\rangle.
\]
\end{remark}

% ---------------------------------------------------------
\subsection{A.9. Нормы тестов шнека: ось e\_\parallel и разведение ``V=0''}

\begin{postulate}[Канонический выбор оси \(e_{\parallel}\)]
Для тестов шнека вводится ось объекта \(e_{\parallel}\) — \textbf{ось симметрии} тора/двутория,
вдоль которой определяется продольный дрейф.

\textbf{Нормативно:} \(e_{\parallel}\) должен быть однозначно воспроизводим из геометрии объекта и его паспорта,
без апелляции к внешним ``полям''.
Для тора/двутория ориентация \(e_{\parallel}\) согласуется с принятой параметризацией \(X(\varphi,\theta)\) поверхности и паспортом вихря:
касательное направление \(\partial_{\varphi}X\) соответствует циркуляции по большой окружности, а \(e_{\parallel}\) — оси симметрии выбранной системы координат.
\end{postulate}

\begin{postulate}[Разведение норм ``\(V=0\)'' и нормы по \(V_{\parallel}\)]
В тестах используются две разные нормы ``нуля'', их нельзя смешивать:
\begin{enumerate}
\item \textbf{Охранная норма (T2/T6):} запись ``\(V=0\)'' означает \textbf{нулевой вектор} скорости дрейфа,
т.е. ноль всех компонент:
\[
V=0 \iff V_x=V_y=V_z=0.
\]
\item \textbf{Норма шнека (T5--T6):} нормативной величиной является \textbf{продольная компонента}
\[
V_{\parallel}\equiv V\cdot e_{\parallel}.
\]
При прочих равных \(\chi\)-нечётной обязана быть именно \(V_{\parallel}\):
\[
V_{\parallel}(\chi)=-V_{\parallel}(-\chi).
\]
При фиксированном \(\chi\) реверс направления бега шнека обязан реверсировать знак \(V_{\parallel}\).
Для стоячего шнека нормативно \(V_{\parallel}=0\) (при отсутствии внешнего нарушителя).
\end{enumerate}
\end{postulate}

\begin{remark}
Смысл нормы: \(\chi\) должна входить в замыкания так, чтобы \(\chi\)-нечётная часть отклика давала вклад в \(V_{\parallel}\),
а не подменялась ``скрытым потоком \SE'' или асимметрией числа контактов.
\end{remark}

\section{Аксиоматическая рамка текущего шага: локальный патч как временное сужение задачи}
\label{sec:a_current_step}

\subsection{Глобальная цель и смысл временного сужения}

Цель текущего шага не состоит в замене глобальной тороидальной или двуториальной
геометрии ЭВО локальным цилиндрическим описанием.
Глобальная цель остаётся прежней:
выявить кинетический механизм дрейфа и ротации ЭВО через контактную статистику
и анизотропию приращений импульса на граничной поверхности $\Sigma$.

Квазицилиндрическое приближение вводится только как временное локальное сужение задачи.
Оно используется там, где режим тонкого тора $r \ll R$ позволяет считать короткий участок
большой окружности локально цилиндрическим.
Тем самым локальный патч $\Sigma_{\mathrm{loc}} \subset \Sigma$
служит не новой онтологией, а увеличительным языком записи,
удобным для выделения локальных механизмов на ГП.

\subsection{Что именно сохраняется неизменным}

В аксиоматическом смысле сохраняются следующие положения:

\begin{enumerate}
\item ЭВО глобально остаётся тороидальной или двуториальной структурой.
\item КУБК относится только к балансу числа контактов на $\Sigma$.
\item Дрейф и вращение допускаются только через анизотропию $\Delta p$ на $\Sigma$.
\item Шнек остаётся динамическим режимом ГП, а не отдельной сущностью и не потоком СЭ.
\item Локальный патч не заменяет глобальный интеграл по $\Sigma$, а только помогает его разложить.
\end{enumerate}

Следовательно, переход к $\Sigma_{\mathrm{loc}}$ не является сменой физической картины,
а является временным усилением разрешающей способности математического языка.

\subsection{Почему непрерывная q-схема оказалась переходной}

На предыдущем этапе непрерывный параметр $q$ использовался как язык локальной намотки.
Однако такой язык оказался недостаточным для сохранения физики шнекового процесса,
поскольку не различал явно:

\begin{itemize}
\item траектории Амеров в канале;
\item шнековую моду граничной поверхности;
\item дискретную связность шнековой организации с канальной навивкой.
\end{itemize}

Поэтому в текущем шаге физический приоритет непрерывного $q$ снимается.
Если параметр $q$ и сохраняется, то только как вспомогательный эффективный
локальный параметр записи $q_{\mathrm{eff}}$.

\subsection{Новая кандидатная рамка текущего шага}

В рабочем контуре проекта принимается кандидатная схема:

\[
m_A \in \mathbb{Z}_{>0},
\qquad
m_S \in \{1,2,4\},
\]

где $m_A$ задаёт квантовое число канальной навивки,
а $m_S$ задаёт базовую пространственную моду шнека на ГП.

Базовыми режимами считаются:
\[
S1,\quad S2,\quad S4,
\]
а все прочие режимы допускаются только как производные от них.

В локальном патч-языке должны различаться два касательных направления:
\[
h_A \quad \text{и} \quad h_S,
\]
где $h_A$ есть директор траектории Амеров,
а $h_S$ --- директор шнековой моды ГП.
Шнек не отождествляется с направлением канальной траектории.

\subsection{Интерпретационная граница}

Текущий шаг проекта следует понимать как последовательность:

\[
\text{глобальная цель}
\;\to\;
\text{тонкий тор}
\;\to\;
\Sigma_{\mathrm{loc}}
\;\to\;
(h_A,h_S)
\;\to\;
(m_A,m_S,S1/S2/S4)
\;\to\;
\text{возврат к глобальному интегралу по } \Sigma.
\]

Тем самым локальный акцент на квазицилиндрическом патче не замещает канон,
а служит промежуточным этапом прояснения механизма.
Именно в таком виде данный шаг должен отражаться в витрине и в журнале изменений.

% Совместимость: article, pdfLaTeX, T2A, babel (ru/uk)
% =========================================================

\section{Аксиоматика и определения A (расширение канона)}

\subsection*{Статус}
\textbf{Ответственный:} Сотрудник A (аксиоматика) \quad
\textbf{Статус:} рабочая вставка \quad
\textbf{Версия:} A-0.2 \quad
\textbf{Дата (Europe/Kyiv):} \today

% ---------------------------------------------------------
\subsection{A.1. Термины и фазовые состояния}

\begin{definition}[Фазовые состояния эфира]
В модели различаются два фазовых состояния движения Амеров:
\begin{itemize}
\item \textbf{Хаотическая фаза:} статистически хаотическое движение Амеров.
Вне \EVO\ хаотическая фаза называется \SE\ (Свободный эфир), а внутри \EVO\ — \NE\ (Несвободный эфир),
включая центральное ядро \NY\ и межканальные прослойки.
\item \textbf{Направленная фаза:} упорядоченное циркуляционное движение, существующее \emph{только} в каналах \EVO.
Направленное движение в каналах \textbf{не} называется \NE.
\end{itemize}
\end{definition}

\begin{axiom}[Запрет устойчивых потоков \SE]
В \SE\ не допускаются устойчивые потоки как первичная онтология.
Допускаются только распространяющиеся \VR\ как флуктуации импульса/плотности,
передаваемые цепочками парных соударений.
\end{axiom}

% ---------------------------------------------------------
\subsection{A.2. Геометрия \EVO: каналы и граничные поверхности}

\begin{definition}[Каналы и граничные поверхности (\GP)]
\EVO\ состоит из одного или нескольких коаксиальных циркуляционных каналов (направленная фаза)
и областей \NE\ (хаос) между ними и в \NY.
Каждый канал имеет две \GP: внутреннюю и внешнюю.
Внешняя \GP внешнего канала является внешней \GP всего \EVO\ (границей \EVO\ с \SE).
\end{definition}

\begin{remark}[Шнековая организация \GP]
Шнекоподобная форма \GP понимается как согласованная пространственно-временная организация контактов,
которая может быть бегущей или стоячей (см. A.6) и влияет на статистику соударений на \GP.
\end{remark}

% ---------------------------------------------------------
\subsection{A.3. Инвариант и жёсткие запреты}

\begin{axiom}[Единственный физический акт]
Единственный физический акт модели — парные упругие соударения Амеров.
Все процессы (структурирование, перенос \VR, взаимодействия, движение/вращение \EVO) реализуются только через них.
\end{axiom}

\begin{axiom}[Жёсткие запреты модели]
Запрещается вводить как первичные сущности:
\begin{itemize}
\item устойчивые потоки \SE;
\item ``давление \SE'' как первичный агент;
\item проникновение Амеров через \GP;
\item поглощение/переизлучение Амеров;
\item ``виртуальные переносчики'' как физические объекты;
\item ``поле'' как отдельную субстанцию (термин допустим только как макро-ярлык).
\end{itemize}
Генезис вихрей из \SE\ в рамках канона не рассматривается (\EVO\ считаются существующими).
\end{axiom}

% ---------------------------------------------------------
\subsection{A.4. Условие целостности хаотической фазы (запрет на обгон)}

\begin{postulate}[Целостность хаотической фазы у \GP]
В хаотической фазе (\SE\ и \NE) у \GP поддерживается режим, обеспечивающий целостность \EVO:
\[
l_{\max}\le d_A, \qquad \mean{l}=0.5\,d_A,
\]
где \(l\) — длина свободного пробега Амера, \(d_A\) — диаметр Амера.
При нарушении режима целостности вихревая структура теряет устойчивость и распадается.
\end{postulate}

\begin{remark}[Каналы]
В каналах (направленная фаза) концентрация ниже, скорости выше;
обгоны/отставания возможны внутри потока, но отсутствует движение против общего направления.
\end{remark}

% ---------------------------------------------------------
\subsection{A.5. КУБК: контактный баланс и минимальная концентрация канала}

\begin{postulate}[КУБК: баланс \emph{числа} контактов на \GP]
В квазистационарном режиме на каждом участке \GP отсутствует устойчивый перекос по \emph{числу} контактов
между амерами хаотической фазы (\SE/\NE) и амерами канала у \GP.
Эквивалентно, после осреднения по времени выполняется согласование \emph{интенсивностей контактов (касания \GP)}:
\[
\Phi^-_{\mathrm{хаос}}(x)=\Phi^+_{\mathrm{канал}}(x),\qquad x\in\Sigma,
\]
где \(\Sigma\) — соответствующая \GP, \(n\) — внешняя единичная нормаль к \(\Sigma\),
а \(\Phi\) не является ``потоком \SE\ в объёме'', а есть локальная величина \emph{частоты/интенсивности касаний поверхности}:
\[
\Phi^-_{\mathrm{хаос}}(x)=\int_{v\cdot n<0}\abs{v\cdot n}\;f_{\mathrm{хаос}}(x,v)\,d^3v,\qquad
\Phi^+_{\mathrm{канал}}(x)=\int_{v\cdot n>0}(v\cdot n)\;f_{\mathrm{канал}}(x,v)\,d^3v.
\]
Здесь \(f_{\mathrm{хаос}}(x,v)\) и \(f_{\mathrm{канал}}(x,v)\) — локальные функции распределения амеров по скоростям
в окрестности \(\Sigma\) со своей стороны поверхности; интегрирование ведётся по соответствующим полупространствам
\(v\cdot n<0\) и \(v\cdot n>0\).
\end{postulate}

\begin{postulate}[Принцип минимальной концентрации канала у \GP]
Среди всех режимов канала, совместимых с целостностью \EVO, реализуется режим с \emph{минимально возможной}
локальной концентрацией \(n_{\mathrm{канал}}(x)\) у \GP, при котором сохраняется КУБК.
\end{postulate}

\begin{remark}[Различение балансов]
КУБК фиксирует баланс \emph{числа} контактов.
Дрейф/вращение определяются не этим балансом, а анизотропией \emph{приращений импульса} в контактах на \GP.
\end{remark}

% ---------------------------------------------------------
\subsection{A.6. Динамическая природа ``шнека''}

\begin{postulate}[Шнек как колебательный режим \GP]
Равновесие на \GP поддерживается колебательным процессом границы относительно среднего положения.
Колебания \GP локально ортогональны направлению основного потока в канале.
Период/фаза колебаний связаны (или фазно согласованы) с характерным временем циркуляции в структуре канала,
в частности с временем обращения вокруг малого кольца тора и/или его гармониками.
Возможны режимы: стоячий шнек и бегущий шнек (по ходу или против хода канального движения),
в зависимости от распределения внутренних/межканальных \VR\ и состояния \KP.
\end{postulate}

% ---------------------------------------------------------
\subsection{A.7. Симметрия дрейфа, хиральность и двуторий}

\begin{postulate}[Симметрия дрейфа в изотропном фоне]
В изотропном фоне (\(\nabla\KP=0\)) и при отсутствии внешних \VR, выделяющих направление,
идеализированный симметричный ``тонкий'' тор не имеет выделенного направления для устойчивого дрейфа.
Следовательно, дрейф требует нарушения симметрии статистики \emph{приращений импульса} на \GP
(внешнего или самосогласованного внутреннего).
\end{postulate}

\begin{definition}[Хиральность (зарученность)]
Хиральность \(\chi\) — дискретный паспортный знак зарученности циркуляции канала/объекта:
\[
\chi\in\{-1,+1\}.
\]
\(\chi\) используется как метка конфигурации (паспорт вихря) и не вводится как отдельная физическая субстанция.
\end{definition}

\begin{definition}[Двуторий]
Двуторий — тороидально-спиральная форма \EVO, возникающая как самосогласованная трансформация ``тонкого'' вихря,
когда собственные \VR\ и связанная с ними структура \KP\ внутри \EVO\ приводят к устойчивому нарушению симметрии
статистики контактов/приращений импульса на \GP.
\end{definition}

\begin{postulate}[Двуторий: механизм ``вращение \(\to\) дрейф'']
В двутории внутренняя топология и шнековая организация \GP обеспечивают преобразование вращательных мод
в поступательный дрейф вдоль оси симметрии.
Механизм трактуется кинетически: интерференция собственных \VR\ формирует устойчивую анизотропию состояния хаотической фазы
(градиенты \KP) у разных участков \GP; интеграл микросмещений мест соударений на \GP даёт ненулевой дрейф и/или момент.
\end{postulate}

\begin{remark}
Формализация дрейфа/момента через поверхностные интегралы от анизотропии импульсного обмена \(j(x,t)\)
является предметом раздела B/Математика.
\end{remark}

% ---------------------------------------------------------
\subsection{A.8. Прокси (K, G, A\_ij): измерительные ярлыки (норматив)}

\begin{axiom}[Прокси как язык измерений, не сущности]
В целях сопоставимости расчётов, симуляций и эксперимента допускается использование прокси-наборов
\((K,\,G,\,A_{ij})\) как \emph{измерительных ярлыков} для состояния хаотической фазы в окрестности \(\Sigma\)
(т.е. для канонических понятий \(\KP\) и \(\VR\)).

\textbf{Нормативно:} прокси не являются новыми сущностями и не добавляют причинности в модель.
Причинность остаётся кинетической: соударения, КУБК (баланс числа контактов) и анизотропия приращений импульса на \(\Sigma\).
\end{axiom}

\begin{remark}[Запретная лексика для прокси]
Запрещены формулировки причинности вида: ``\(G\) действует/толкает'',
``тензор \(A_{ij}\) создаёт силу'', ``поле \(K\) ускоряет'' и т.п.
Допустимо только: ``прокси \(G\)/\(A_{ij}\) наблюдается/оценён'' и ``используется как измерительная метка''.
\end{remark}

\begin{remark}[Рекомендуемая нормировка (не обязательна)]
Допускается использовать единый вид для первых моментов распределения:
\[
K=n\langle v^2\rangle,\qquad G=n\langle v\rangle,\qquad
A_{ij}=P_{ij}-\frac{1}{3}\tr(P)\,\delta_{ij},\qquad P_{ij}=n\langle v_i v_j\rangle.
\]
\end{remark}

% ---------------------------------------------------------
\subsection{A.9. Нормы тестов шнека: ось e\_\parallel и разведение ``V=0''}

\begin{postulate}[Канонический выбор оси \(e_{\parallel}\)]
Для тестов шнека вводится ось объекта \(e_{\parallel}\) — \textbf{ось симметрии} тора/двутория,
вдоль которой определяется продольный дрейф.

\textbf{Нормативно:} \(e_{\parallel}\) должен быть однозначно воспроизводим из геометрии объекта и его паспорта,
без апелляции к внешним ``полям''.
Для тора/двутория ориентация \(e_{\parallel}\) согласуется с принятой параметризацией \(X(\varphi,\theta)\) поверхности и паспортом вихря:
касательное направление \(\partial_{\varphi}X\) соответствует циркуляции по большой окружности, а \(e_{\parallel}\) — оси симметрии выбранной системы координат.
\end{postulate}

\begin{postulate}[Разведение норм ``\(V=0\)'' и нормы по \(V_{\parallel}\)]
В тестах используются две разные нормы ``нуля'', их нельзя смешивать:
\begin{enumerate}
\item \textbf{Охранная норма (T2/T6):} запись ``\(V=0\)'' означает \textbf{нулевой вектор} скорости дрейфа,
т.е. ноль всех компонент:
\[
V=0 \iff V_x=V_y=V_z=0.
\]
\item \textbf{Норма шнека (T5--T6):} нормативной величиной является \textbf{продольная компонента}
\[
V_{\parallel}\equiv V\cdot e_{\parallel}.
\]
При прочих равных \(\chi\)-нечётной обязана быть именно \(V_{\parallel}\):
\[
V_{\parallel}(\chi)=-V_{\parallel}(-\chi).
\]
При фиксированном \(\chi\) реверс направления бега шнека обязан реверсировать знак \(V_{\parallel}\).
Для стоячего шнека нормативно \(V_{\parallel}=0\) (при отсутствии внешнего нарушителя).
\end{enumerate}
\end{postulate}

\begin{remark}
Смысл нормы: \(\chi\) должна входить в замыкания так, чтобы \(\chi\)-нечётная часть отклика давала вклад в \(V_{\parallel}\),
а не подменялась ``скрытым потоком \SE'' или асимметрией числа контактов.
\end{remark}

\section{Аксиоматическая рамка текущего шага: локальный патч как временное сужение задачи}
\label{sec:a_current_step}

\subsection{Глобальная цель и смысл временного сужения}

Цель текущего шага не состоит в замене глобальной тороидальной или двуториальной
геометрии ЭВО локальным цилиндрическим описанием.
Глобальная цель остаётся прежней:
выявить кинетический механизм дрейфа и ротации ЭВО через контактную статистику
и анизотропию приращений импульса на граничной поверхности $\Sigma$.

Квазицилиндрическое приближение вводится только как временное локальное сужение задачи.
Оно используется там, где режим тонкого тора $r \ll R$ позволяет считать короткий участок
большой окружности локально цилиндрическим.
Тем самым локальный патч $\Sigma_{\mathrm{loc}} \subset \Sigma$
служит не новой онтологией, а увеличительным языком записи,
удобным для выделения локальных механизмов на ГП.

\subsection{Что именно сохраняется неизменным}

В аксиоматическом смысле сохраняются следующие положения:

\begin{enumerate}
\item ЭВО глобально остаётся тороидальной или двуториальной структурой.
\item КУБК относится только к балансу числа контактов на $\Sigma$.
\item Дрейф и вращение допускаются только через анизотропию $\Delta p$ на $\Sigma$.
\item Шнек остаётся динамическим режимом ГП, а не отдельной сущностью и не потоком СЭ.
\item Локальный патч не заменяет глобальный интеграл по $\Sigma$, а только помогает его разложить.
\end{enumerate}

Следовательно, переход к $\Sigma_{\mathrm{loc}}$ не является сменой физической картины,
а является временным усилением разрешающей способности математического языка.

\subsection{Почему непрерывная q-схема оказалась переходной}

На предыдущем этапе непрерывный параметр $q$ использовался как язык локальной намотки.
Однако такой язык оказался недостаточным для сохранения физики шнекового процесса,
поскольку не различал явно:

\begin{itemize}
\item траектории Амеров в канале;
\item шнековую моду граничной поверхности;
\item дискретную связность шнековой организации с канальной навивкой.
\end{itemize}

Поэтому в текущем шаге физический приоритет непрерывного $q$ снимается.
Если параметр $q$ и сохраняется, то только как вспомогательный эффективный
локальный параметр записи $q_{\mathrm{eff}}$.

\subsection{Новая кандидатная рамка текущего шага}

В рабочем контуре проекта принимается кандидатная схема:

\[
m_A \in \mathbb{Z}_{>0},
\qquad
m_S \in \{1,2,4\},
\]

где $m_A$ задаёт квантовое число канальной навивки,
а $m_S$ задаёт базовую пространственную моду шнека на ГП.

Базовыми режимами считаются:
\[
S1,\quad S2,\quad S4,
\]
а все прочие режимы допускаются только как производные от них.

В локальном патч-языке должны различаться два касательных направления:
\[
h_A \quad \text{и} \quad h_S,
\]
где $h_A$ есть директор траектории Амеров,
а $h_S$ --- директор шнековой моды ГП.
Шнек не отождествляется с направлением канальной траектории.

\subsection{Интерпретационная граница}

Текущий шаг проекта следует понимать как последовательность:

\[
\text{глобальная цель}
\;\to\;
\text{тонкий тор}
\;\to\;
\Sigma_{\mathrm{loc}}
\;\to\;
(h_A,h_S)
\;\to\;
(m_A,m_S,S1/S2/S4)
\;\to\;
\text{возврат к глобальному интегралу по } \Sigma.
\]

Тем самым локальный акцент на квазицилиндрическом патче не замещает канон,
а служит промежуточным этапом прояснения механизма.
Именно в таком виде данный шаг должен отражаться в витрине и в журнале изменений.


% =========================================================
\section{Математическое ядро  B}
\subsection*{Статус}
\textbf{Версия:} v0.2 \quad
\textbf{Дата (Europe/Kyiv):} \today \quad
\textbf{Ответственный:} Сотрудник B (математика)
\label{sec:math-core}
% ВАЖНО: sec_kinetics_B лежит рядом с main.tex
% =========================
% sec_kinetics_B.tex — ПОЛНАЯ ПЕРЕЗАПИСЬ ДЛЯ ВИТРИНЫ (B/Математика)
%
% Витринная логика:
% 1) не подменять глобальную тороидальную геометрию локальным патчем;
% 2) показать, зачем вводится Σ_loc;
% 3) сохранить канон: дрейф/вращение только через анизотропию Δp на Σ;
% 4) показать эволюцию языка: от q-схемы к дискретно-модовой схеме m_A, m_S, S1/S2/S4;
% 5) вернуть локальные формулы к глобальному интегралу по Σ.
%
% Требует в преамбуле:
% \usepackage{tikz}
% \usepackage{booktabs}
% \usepackage{array}
% \usetikzlibrary{arrows.meta,calc,positioning,decorations.pathmorphing}
%
% Канон-опора:
% - CANON.md — единственный источник истины;
% - PDF/LaTeX — витрина и может отставать;
% - тонкий тор r << R допускает только локальное квазицилиндрическое приближение;
% - шнек — динамический режим ГП;
% - КУБК = баланс числа контактов, а не источник дрейфа;
% - дрейф/вращение только через анизотропию Δp на Σ;
% - прокси (K,G,A_ij) — ярлыки измерения;
% - e_∥ — ось §17, не h_A и не h_S.
% =========================

\section{Кинетика ЭВО: дрейф и вращение как интегральный отклик анизотропии контактов на \(\Sigma\)}

\subsection{Глобальная цель и временное сужение задачи}

Цель текущего шага --- не заменить глобальную тороидальную геометрию ЭВО локальным цилиндром, а временно сузить задачу до локального патча \(\Sigma_{\mathrm{loc}}\), чтобы явнее выделить кинетику контактов на ГП, ответственную за дрейф и ротацию. Глобально ЭВО остаётся тороидальной или двуториальной структурой; локальный патч используется только как язык записи и вычисления для той же самой канонической геометрии.

Витринно важно сохранить последовательность рассуждения. Сначала задаётся глобальная геометрия и полный интеграл по \(\Sigma\). Затем объясняется, почему в режиме тонкого тора возникает естественная локальная карта \(\Sigma_{\mathrm{loc}}\). После этого становится видно, какие механизмы удобнее отделяются на патче, почему одной непрерывной \(q\)-геометризации оказалось недостаточно, почему пришлось различить \(h_A\) и \(h_S\), и почему шнек теперь следует читать через дискретно-согласованную схему \((m_A,m_S)\) с базовыми модами \(S1/S2/S4\). Финальный шаг --- возврат к глобальному интегралу по \(\Sigma\), чтобы локальный язык не подменял онтологию объекта.

\subsection{Тонкий тор как базовая геометрия и причины локального патча}

\subsubsection{Граничная поверхность \(\Sigma\)}

Пусть \(\Sigma\subset\mathbb{R}^3\) --- компактная кусочно-\(C^2\) ориентируемая поверхность без края
(или конечное объединение таких компонент для многоканального ЭВО). Задана параметризация
\[
X:U\subset\mathbb{R}^2\to\Sigma,\qquad x=X(u,v).
\]
Элемент площади и единичная нормаль:
\[
dS=\|X_u\times X_v\|\,dudv,\qquad
n(x)=\frac{X_u\times X_v}{\|X_u\times X_v\|}.
\]

\subsubsection{Ось \(e_\parallel\) (норма §17)}

Ось \(e_\parallel\) --- \textbf{ось симметрии объекта} и \textbf{ось тестов §17}:
\[
V_\parallel := V\cdot e_\parallel.
\]
Везде далее \(e_\parallel\) используется как нормативная ``χ-ручка''. Важно: \(e_\parallel\) \emph{не} подменяется ни директором траектории Амеров \(h_A\), ни директором шнековой моды \(h_S\).

\subsubsection{Хиральность \(\chi\)}

Фиксируем паспортную хиральность:
\[
\chi\in\{-1,+1\}.
\]

\subsubsection{Контактная статистика на \(\Sigma\) и КУБК}

На \(\Sigma\) необходимо жёстко различать:
\begin{itemize}
\item баланс \emph{числа контактов} --- КУБК;
\item баланс \emph{приращений импульса} --- источник дрейфа и момента.
\end{itemize}

Пусть \(f_{\mathrm{хаос}}(x,v,t)\) и \(f_{\mathrm{канал}}(x,v,t)\) --- локальные распределения скоростей амеров
в окрестности \(\Sigma\) со стороны хаотической фазы и со стороны канала. Нормаль \(n(x)\) направлена из ЭВО наружу.

Интенсивности контактов:
\[
\Phi^-_{\mathrm{хаос}}(x,t)=\int_{v\cdot n\le 0}|v\cdot n|\;f_{\mathrm{хаос}}(x,v,t)\,d^3v,\qquad
\Phi^+_{\mathrm{канал}}(x,t)=\int_{v\cdot n\ge 0}(v\cdot n)\;f_{\mathrm{канал}}(x,v,t)\,d^3v.
\]

\(\Phi^\pm\) трактуются только как \textbf{интенсивности касаний \(\Sigma\)}, а не как ``потоки СЭ''.

КУБК фиксирует отсутствие устойчивого перекоса по \emph{числу} контактов:
\[
\overline{\Phi^-_{\mathrm{хаос}}}(x)=\overline{\Phi^+_{\mathrm{канал}}}(x),\qquad x\in\Sigma.
\]

Это означает: устойчивый дрейф нельзя выводить из перекоса числа контактов. Он может возникать только через анизотропию \(\Delta p\) на \(\Sigma\).

\subsubsection{Полный интеграл по \(\Sigma\)}

Локальная импульсная анизотропия задаётся полем
\[
j(x,t):=\lambda(x,t)\,\mathbb{E}[\Delta p\mid x,t],\qquad x\in\Sigma.
\]
Разложение:
\[
j=j_n\,n+j_\top,\qquad j_n=(j\cdot n),\qquad j_\top\perp n.
\]

Интегральные отклики:
\[
V(t)=\mu_T\int_{\Sigma} j(x,t)\,dS,
\]
\[
\Omega(t)=\tilde\mu_R\int_{\Sigma}\bigl(r(x)-r_0\bigr)\times j(x,t)\,dS.
\]

Именно этот глобальный интеграл по \(\Sigma\) остаётся главным объектом кинетики. Все локальные записи ниже должны возвращаться к нему.

\subsection{Квазицилиндрический патч как локальный язык, а не новая онтология}

\subsubsection{Почему нужен локальный патч}

В режиме тонкого тора \(r\ll R\) изменение локальной кривизны по большой окружности мало на масштабе длины свободного пробега \(l\). Поэтому короткий участок \(\Sigma\) можно записывать как квазицилиндрический патч:
\[
\Sigma_{\mathrm{loc}}\subset\Sigma.
\]

Это не новая геометрия объекта и не альтернатива тору. Это локальная карта, в которой проще отделить:
\begin{itemize}
\item контактную статистику на ГП;
\item касательные компоненты \(j(x)\);
\item локальные tor/sp-вклады;
\item различие между траекторией Амеров и шнековой модой ГП.
\end{itemize}

\subsubsection{Локальный базис патча}

На \(\Sigma_{\mathrm{loc}}\) вводим:
\[
(e_z,e_\vartheta,n),
\]
где
\begin{itemize}
\item \(e_z\) --- единичный касательный вектор вдоль большой окружности тора;
\item \(e_\vartheta\) --- единичный касательный вектор по малой окружности;
\item \(n\) --- внешняя единичная нормаль.
\end{itemize}

Локальное соответствие глобальной тороидальной геометрии:
\[
e_z \leftrightarrow t_\varphi,\qquad e_\vartheta \leftrightarrow t_\theta.
\]

Таким образом, патч --- это лишь локальный язык записи той же граничной поверхности \(\Sigma\).

\subsubsection{Прокси как язык наблюдения}

Сохраняем канонический прокси-стандарт:
\[
K=n\langle v^2\rangle,\qquad
G=n\langle v\rangle,\qquad
P_{ij}=n\langle v_i v_j\rangle,\qquad
A_{ij}=P_{ij}-\frac13\mathrm{tr}(P)\delta_{ij}.
\]
Проекции на касательную плоскость:
\[
\Pi_\top(\xi)=\xi-(\xi\cdot n)n,\qquad
g_\top:=\Pi_\top(G),\qquad
a_\top:=\Pi_\top(A\,n).
\]

Прокси не являются сущностями и не несут причинности; они служат только языком измерения контактной статистики и анизотропии \(\Delta p\).

% =========================================================
% ВСТАВКА 1: БОЛЬШАЯ ВИЗУАЛЬНАЯ СХЕМА
% =========================================================

\subsubsection{Схема локального квазицилиндрического патча и направлений кинетики}

На рис.~\ref{fig:quasi_cyl_patch_scheme} показана расчётная схема локального квазицилиндрического приближения. Глобально ЭВО сохраняет тороидальную геометрию с граничной поверхностью \(\Sigma\). В режиме тонкого тора \(r\ll R\) выделяется локальный патч \(\Sigma_{\mathrm{loc}}\), на котором вводятся координаты \((z,\vartheta)\), внешняя нормаль \(n\), касательные направления \(e_z\) и \(e_\vartheta\), а также два различных касательных директора: \(h_A\) для траекторий Амеров в канале и \(h_S\) для шнековой моды ГП. Важно, что поток Амеров и бег шнека не отождествляются: первый связан с \(h_A\), второй --- с \(h_S\). Локально вычисляются контактные интенсивности, прокси и внутренние вклады \(j_{\mathrm{tor}}^{\mathrm{loc}}\), \(j_{\mathrm{sp}}^{\mathrm{loc}}\), после чего результат собирается обратно в глобальный интеграл по \(\Sigma\).

% =========================================================
% ПОЛНОСТРАНИЧНАЯ ВСТАВКА:
% крупная схема квазицилиндрического патча
%
% Требует в преамбуле:
% \usepackage{tikz}
% \usepackage{graphicx}
% \usetikzlibrary{arrows.meta,calc,positioning,decorations.pathmorphing}
% =========================================================

% =========================================================
% ПОЛНОСТРАНИЧНАЯ ВЕРТИКАЛЬНАЯ СХЕМА:
% рисунки расположены один над другим
%
% Требует в преамбуле:
% \usepackage{tikz}
% \usepackage{graphicx}
% \usetikzlibrary{arrows.meta,calc,positioning,decorations.pathmorphing}
% =========================================================

% =========================================================
% ПОЛНОСТРАНИЧНАЯ ВЕРТИКАЛЬНАЯ СХЕМА
% Готовый блок с уменьшенными шрифтами
% =========================================================

% =========================================================
% ПОЛНОСТРАНИЧНАЯ ВЕРТИКАЛЬНАЯ СХЕМА
% Готовый блок с уменьшенными шрифтами
% =========================================================

% =========================================================
% ОТДЕЛЬНЫЙ ЛИСТ С РИСУНКОМ НА ВСЮ ПЛОЩАДЬ
% Без caption, чтобы схема заняла максимум места
% =========================================================

% =========================================================
% ПОЛНОСТРАНИЧНАЯ СХЕМА — СИЛЬНО ОБЛЕГЧЁННАЯ
% Внутри рисунка только короткие метки
% =========================================================

% =========================================================
% НОВАЯ СХЕМА С НУЛЯ
% Полностраничный рисунок, минимальный текст внутри
%
% Требует в преамбуле:
% \usepackage{tikz}
% \usepackage{graphicx}
% \usepackage{geometry}
% \usetikzlibrary{arrows.meta,calc,positioning}
% =========================================================

% =========================================================
% РИСУНОК А — ГЕОМЕТРИЯ КВАЗИЦИЛИНДРИЧЕСКОГО ПАТЧА
% Полностраничный, без перегруза формулами
%
% Требует в преамбуле:
% \usepackage{tikz}
% \usepackage{geometry}
% \usetikzlibrary{arrows.meta,calc,positioning}
% =========================================================

\clearpage
\begingroup
\thispagestyle{empty}
\newgeometry{top=5mm,bottom=5mm,left=5mm,right=5mm}

\phantomsection
\refstepcounter{figure}
\label{fig:quasi_cyl_geometry}
\addcontentsline{lof}{figure}{Геометрия локального квазицилиндрического патча}

\begin{center}
\resizebox{0.99\textwidth}{0.96\textheight}{%
\begin{tikzpicture}[
    >=Latex,
    font=\Large,
    boundary/.style={line width=1.3pt},
    axis/.style={->, line width=1.25pt},
    amer/.style={->, line width=1.3pt},
    screw/.style={->, dashed, line width=1.3pt},
    lbl/.style={font=\Large},
    sml/.style={font=\large},
    every node/.style={inner sep=2pt}
]

% ---------------------------------------------------------
% Верх: глобальный тор
% ---------------------------------------------------------

\node[lbl] at (-3.8,11.1) {\(\Sigma\)};

\begin{scope}[shift={(0,8.25)}]
    \draw[boundary] (0,0) ellipse (3.35 and 2.0);
    \draw[boundary] (0,0) ellipse (1.58 and 0.95);

    \draw[line width=1.8pt] (-1.15,0.72)
        arc[start angle=145,end angle=35,x radius=1.58,y radius=0.95];

    \draw[rounded corners=5pt, line width=1.2pt] (0.95,0.38) rectangle (2.45,1.28);
    \node[sml] at (3.55,1.15) {\(\Sigma_{\mathrm{loc}}\)};

    \draw[axis] (0,0) -- (2.8,0) node[midway, below] {\(R\)};
    \draw[axis] (0,0) -- (0,1.05) node[midway, right] {\(r\)};
\end{scope}

\draw[axis] (5.0,8.3) -- (6.65,8.3);
\node[lbl] at (5.8,9.0) {\(r\ll R\)};

% ---------------------------------------------------------
% Середина: локальный патч
% ---------------------------------------------------------

\node[lbl] at (0,6.15) {\(\Sigma_{\mathrm{loc}}\)};

\begin{scope}[shift={(-4.0,0.95)}]
    % рамка патча
    \draw[boundary] (0,0) rectangle (8.45,5.35);

    % координаты
    \node[sml] at (4.2,-0.55) {\(z\)};
    \node[sml, rotate=90] at (-0.7,2.65) {\(\vartheta\)};

    % базис
    \draw[axis] (0.95,0.85) -- (2.95,0.85) node[right] {\(e_z\)};
    \draw[axis] (0.95,0.85) -- (0.95,2.85) node[above] {\(e_\vartheta\)};

    % точка и нормаль
    \fill (3.2,2.55) circle (2pt);
    \node[sml, above right] at (3.2,2.55) {\(x\)};
    \draw[axis] (3.2,2.55) -- (4.5,3.95) node[above] {\(n(x)\)};

    % вертикальные стрелки контактов сверху
    \foreach \xx in {1.0,1.9,3.0,4.15,5.45,6.75} {
        \draw[axis] (\xx,6.05) -- (\xx,5.35);
    }
    \node[sml] at (2.75,6.45) {контакты};

    % шнековые линии
    \foreach \c in {-0.8,0.15,1.1,2.05} {
        \draw[screw] (0.45,\c+0.95) -- (4.25,\c+4.05);
    }
    \node[sml] at (5.95,4.25) {\(s_{\rm sh}\)};
    \draw[screw] (5.0,3.8) -- (5.85,4.35);

    % траектории Амеров
    \foreach \yy in {1.25,1.95,2.65,3.35,4.05} {
        \draw[amer] (5.1,\yy) -- (7.15,\yy+0.46);
    }
    \node[sml] at (7.25,2.65) {Амеры};

    % h_A и h_S
    \draw[amer] (3.2,2.55) -- (4.9,2.95);
    \node[sml, below right] at (4.9,2.95) {\(h_A\)};

    \draw[screw] (3.2,2.55) -- (4.65,4.0);
    \node[sml, above right] at (4.65,4.0) {\(h_S\)};

    % краткие поля
    \node[sml, align=left] at (1.95,3.55) {\(\Phi^\pm,\ \lambda\)};
    \node[sml, align=left] at (2.0,2.95) {\(K,\ G,\ A_{ij}\)};
    \node[sml, align=left] at (2.25,1.25) {\(j_{\mathrm{tor}}^{\mathrm{loc}},\ j_{\mathrm{sp}}^{\mathrm{loc}}\)};

    % различение директоров
    \node[sml, align=left] at (6.55,1.0) {\(h_A \neq h_S\)};
    \node[sml, align=left] at (6.65,0.45) {\(e_\parallel \neq h_A,h_S\)};
\end{scope}

% ---------------------------------------------------------
% Низ: минимальный вывод из геометрии
% ---------------------------------------------------------

\node[lbl] at (0,-1.15) {\(\Sigma_{\mathrm{loc}}\subset\Sigma\)};
\node[lbl] at (0,-2.35) {\((e_z,e_\vartheta,n),\ h_A,\ h_S\)};
\node[lbl] at (0,-3.55) {\(r\ll R\ \Rightarrow\) локальная квазицилиндрическая запись};

\end{tikzpicture}%
}
\end{center}

\restoregeometry
\endgroup
\clearpage






% =========================================================
% ВСТАВКА 2: КОМПАКТНАЯ ЛЕГЕНДА / ТАБЛИЦА К РИСУНКУ
% =========================================================

\subsubsection{Легенда к локальному патчу и дискретно-модовой записи}

Для удобства чтения соберём основные обозначения, используемые на рис.~\ref{fig:quasi_cyl_patch_scheme} и в локальных формулах, в компактную таблицу.

\begin{table}[htbp]
\centering
\small
\begin{tabular}{>{\raggedright\arraybackslash}p{2.2cm} >{\raggedright\arraybackslash}p{9.8cm}}
\toprule
Обозначение & Смысл \\
\midrule
\(\Sigma\) & глобальная граничная поверхность ЭВО \\
\(\Sigma_{\mathrm{loc}}\) & локальный квазицилиндрический патч на \(\Sigma\) \\
\(z\) & координата вдоль большой окружности тора \\
\(\vartheta\) & координата по малой окружности тора \\
\(n\) & внешняя единичная нормаль к \(\Sigma_{\mathrm{loc}}\) \\
\(e_z\) & касательное направление вдоль большой окружности (\(e_z\leftrightarrow t_\varphi\)) \\
\(e_\vartheta\) & касательное направление по малой окружности (\(e_\vartheta\leftrightarrow t_\theta\)) \\
\(e_\parallel\) & ось симметрии объекта и ось тестов §17; не совпадает с \(h_A\) и \(h_S\) \\
\(h_A\) & директор траекторий Амеров в канале; задаётся канальной навивкой \(m_A\) \\
\(h_S\) & директор шнековой моды ГП; задаётся модой \(m_S\), хиральностью \(\chi\) и направлением бега \(s_{\rm sh}\) \\
\(\Phi^\pm,\ \lambda\) & интенсивности контактов на ГП \\
\(K,G,A_{ij}\) & прокси-ярлыки локальной статистики контактов \\
\(g_\top,\ a_\top\) & касательные проекции \(G\) и \(A n\) \\
\(m_A\) & квантовое число канальной навивки траекторий Амеров \\
\(m_S\) & базовая шнековая мода ГП \\
\(S1,S2,S4\) & базовые моды шнека: 1, 2 и 4 оборота Амера \\
\(S^\*\) & производная/гармоническая/составная мода \\
\(q_{\mathrm{eff}}\) & только эффективный локальный наклон патча; не главный физический носитель шнека \\
\(j_{\mathrm{tor}}^{\mathrm{loc}}\) & локальный тороидальный вклад внутренней кинетики \\
\(j_{\mathrm{sp}}^{\mathrm{loc}}\) & локальный пространственно-дрейфовый вклад, производный от tor-вклада \\
\(j_{VR}^{\mathrm{int,loc}}\) & полный внутренний вклад на патче \\
\bottomrule
\end{tabular}
\caption{
Компактная легенда к локальному квазицилиндрическому патчу. 
Таблица подчёркивает, что \(\Sigma_{\mathrm{loc}}\) является только локальным языком записи и вычисления, а физически ведущими параметрами шнекового режима служат дискретные величины \(m_A\) и \(m_S\), а не непрерывный параметр наклона.
}
\label{tab:quasi_cyl_patch_legend}
\end{table}

Дополнительно удобно явно выписать локальную последовательность вычисления:
\[
j(x)=\lambda(x)\,\mathbb{E}[\Delta p\mid x],
\]
\[
j_{\mathrm{tor}}^{\mathrm{loc}}
=
\chi\, s_{\mathrm{sh}}\,
\mathcal{T}(m_A;K,g_\top,a_\top,n),
\]
\[
j_{\mathrm{sp}}^{\mathrm{loc}}
=
\beta(K;m_A,m_S)\,
w_{\mathrm{loc}}(x)\,
\Gamma_{A,S}(x)\,
h_S(x),
\qquad
\Gamma_{A,S}(x)=
\bigl(j_{\mathrm{tor}}^{\mathrm{loc}}\cdot h_A\bigr)\,
\Xi(m_A,m_S),
\]
\[
j_{VR}^{\mathrm{int,loc}}
=
j_{\mathrm{tor}}^{\mathrm{loc}}
+
j_{\mathrm{sp}}^{\mathrm{loc}},
\]
после чего локальная запись снова собирается в глобальные интегралы по \(\Sigma\).

\subsection{Почему непрерывный параметр \(q\) оказался недостаточным}

Прежняя непрерывная \(q\)-геометризация была полезным промежуточным языком. Она позволила впервые выразить идею, что пространственный дрейф двутория невозможен без зашитой в геометрию винтовой организации, а также развести тороидальную и пространственную компоненты внутреннего вклада.

Исторически это выглядело так:
\[
h(x)=\frac{t_\varphi(x)+q\,t_\theta(x)}{\|t_\varphi(x)+q\,t_\theta(x)\|},
\qquad
q\in\mathbb{R}\setminus\{0\},
\]
и затем пространственный вклад записывался как производный от tor-вклада через этот директор намотки.

Однако такой язык оказался недостаточным по физическому смыслу. Он позволял параметризовать локальный наклон, но не удерживал квантовую связность шнекового процесса с канальной навивкой. Если шнек действительно обусловлен квантовыми числами каналов и навивкой траекторий Амеров, то одного непрерывного \(q\) уже недостаточно, чтобы не потерять физику процесса.

Поэтому в текущей редакции \(q\) снимается как главный физический носитель шнека. Он допускается лишь как вспомогательный эффективный параметр локального наклона:
\[
q_{\mathrm{eff}},
\]
и не используется как самостоятельное основание физической интерпретации режима.

\subsection{Два касательных директора: \(h_A\) и \(h_S\)}

Ключевое уточнение нового каркаса состоит в том, что на \(\Sigma_{\mathrm{loc}}\) нужно различать два касательных семейства.

\subsubsection{Директор траектории Амеров}

\[
h_A=h_A(m_A;x),
\]
где \(m_A\in\mathbb{Z}_{>0}\) --- дискретное квантовое число канальной навивки. Вектор \(h_A\) задаёт локальное касательное направление, связанное именно с траекторией Амеров в канале.

\subsubsection{Директор шнековой моды ГП}

\[
h_S=h_S(m_S,\chi,s_{\mathrm{sh}};x),
\]
где \(m_S\) задаёт дискретную базовую пространственную моду шнека на ГП. Вектор \(h_S\) описывает не канал как таковой, а шнековую организацию самой граничной поверхности.

\subsubsection{Почему это разведение необходимо}

Шнек не отождествляется с канальной траекторией. Если их не развести, то шнековая геометрия незаметно растворяется в одной непрерывной ``наклонной'' записи и перестаёт быть отличимой как отдельный динамический режим ГП.

Именно поэтому в новом каркасе:
\begin{itemize}
\item \(h_A\) относится к канальной навивке;
\item \(h_S\) относится к шнековой моде ГП;
\item \(e_\parallel\) остаётся осью §17 и не совпадает ни с \(h_A\), ни с \(h_S\).
\end{itemize}

Если на патче используется гипотеза двух винтовых семейств, она должна быть явно помечена как кандидат:
\[
h_A\cdot h_S=0.
\]
При необходимости допускается и кандидат-гипотеза о специальном угле амерной навивки на развёртке, но она не должна быть скрыта внутри параметра.

\subsection{Дискретно-модовая схема шнека: \(m_A\), \(m_S\), \(S1/S2/S4\)}

\subsubsection{Почему шнек теперь читается дискретно}

Теперь физика шнека читается не через произвольный непрерывный наклон, а через согласованную пару
\[
(m_A,m_S),
\]
где:
\[
m_A\in\mathbb{Z}_{>0}
\]
--- квантовое число канальной навивки траекторий Амеров,
а
\[
m_S\in\{1,2,4\}
\]
--- базовая пространственная мода шнека на ГП.

Это выражает новую рабочую норму: шнековая навивка не свободна, а согласована с канальной навивкой.

\subsubsection{Таблица базовых и производных режимов}

\begin{center}
\begin{tabular}{|c|c|c|p{7.2cm}|}
\hline
Обозначение & Параметр & Статус & Смысл \\
\hline
\(S1\) & \(m_S=1\) & базовая мода & шнек растянут на 1 оборот Амера \\
\hline
\(S2\) & \(m_S=2\) & базовая мода & шнек растянут на 2 оборота Амера \\
\hline
\(S4\) & \(m_S=4\) & базовая мода & шнек растянут на 4 оборота Амера \\
\hline
\(S^\*\) & \(m_S\notin\{1,2,4\}\) & только производная & гармоническая/составная мода, допустимая только как производная от \(S1/S2/S4\) \\
\hline
\end{tabular}
\end{center}

Таблица нужна не как голая классификация, а как ответ на физический вопрос: почему именно эти режимы следует считать базовыми. В текущем каркасе это те моды, которые принимаются как первичные пространственные режимы шнековой организации ГП; все остальные допустимы лишь как производные и не должны интерпретироваться как новые базовые формы без отдельного обоснования.

\subsubsection{Локальные вклады в новой записи}

Локальный tor-вклад:
\[
j_{\mathrm{tor}}^{\mathrm{loc}}
=
\chi\, s_{\mathrm{sh}}\,
\mathcal{T}(m_A;K,g_\top,a_\top,n),
\]
где минимальная рабочая запись:
\[
\mathcal{T}(m_A;K,g_\top,a_\top,n)
=
\alpha_1(K,m_A)\,n\times g_\top
+
\alpha_2(K,m_A)\,n\times a_\top.
\]

Локальный sp-вклад:
\[
j_{\mathrm{sp}}^{\mathrm{loc}}
=
\beta(K;m_A,m_S)\,
w_{\mathrm{loc}}(x)\,
\Gamma_{A,S}(x)\,
h_S(x),
\]
где
\[
\Gamma_{A,S}(x)=
\bigl(j_{\mathrm{tor}}^{\mathrm{loc}}\cdot h_A\bigr)\,
\Xi(m_A,m_S).
\]

Полный внутренний вклад:
\[
j_{VR}^{\mathrm{int,loc}}
=
j_{\mathrm{tor}}^{\mathrm{loc}}
+
j_{\mathrm{sp}}^{\mathrm{loc}}.
\]

В этой записи сохраняется главный принцип: без tor-режима дрейфа нет, а пространственный вклад строго производен от тороидального.

\subsection{Возврат от локального патча к глобальному интегралу по \(\Sigma\)}

Локальный патч не завершает картину, а только подготавливает её. Итоговый физический отклик остаётся глобальным:
\[
V=\mu_T\int_{\Sigma} j\,dS,
\qquad
\Omega=\tilde\mu_R\int_{\Sigma}(r-r_0)\times j\,dS,
\qquad
V_\parallel=V\cdot e_\parallel.
\]

Практически это означает, что глобальный интеграл может быть собран как сумма по локальным патчам:
\[
\int_{\Sigma}(\cdots)\,dS
\approx
\sum_{\mathrm{patches}}
\int_{\Sigma_{\mathrm{loc}}}(\cdots)\,dS.
\]

Но именно здесь нужно поставить последнюю витринную оговорку: локальная запись не имеет собственной онтологии. Она не заменяет тор и не даёт права интерпретировать \(\Sigma_{\mathrm{loc}}\) как самостоятельную физическую форму ЭВО. Смысл локального патча --- лишь в том, чтобы яснее показать механизм, который потом снова собирается в глобальный тороидальный отклик.

\subsection{Охранные нормы и интерпретационные границы}

Для витрины полезно завершить раздел коротким блоком ограничений:

\begin{itemize}
\item КУБК относится только к балансу числа контактов на \(\Sigma\).
\item Дрейф и вращение возникают только через анизотропию \(\Delta p\) на \(\Sigma\).
\item Прокси \((K,G,A_{ij})\) --- только ярлыки измерения.
\item \(q_{\mathrm{eff}}\) не является главным физическим носителем шнека.
\item Базовые моды шнека читаются через \(m_S\in\{1,2,4\}\), а не через произвольную непрерывную деформацию.
\item Различение \(h_A\) и \(h_S\) обязательно.
\item Любая локальная запись на \(\Sigma_{\mathrm{loc}}\) должна возвращаться к глобальному интегралу по \(\Sigma\).
\end{itemize}

\subsection{Переходный статус старой \(q\)-схемы}

Для непрерывности изложения полезно сохранить короткое историческое замечание.

\paragraph{Исторический слой.}
Прежняя непрерывная \(q\)-геометризация была полезным промежуточным языком, потому что позволила отделить чисто тороидальный вклад от пространственно-дрейфового и впервые зафиксировала, что дрейф двутория не может рождаться без геометрически зашитой винтовой организации. Однако в дальнейшем оказалось, что одной непрерывной \(q\)-параметризации недостаточно для удержания физики шнекового процесса, так как она не фиксирует жёсткую связь шнека с канальной навивкой и квантовыми числами каналов.

Поэтому в текущей редакции \(q\)-схема остаётся в витрине как исторический переходный слой, а физически ведущей становится дискретно-модовая схема:
\[
(m_A,m_S),\qquad S1/S2/S4,\qquad h_A,\ h_S,\qquad q_{\mathrm{eff}}\ \text{только как вспомогательный параметр.}
\]

% Конец файла


% =========================================================
\section{Следствия, предсказания и тесты (реестр)}
\subsection*{Статус}
\textbf{Версия:} v0.1 \quad
\textbf{Дата (Europe/Kyiv):} \today \quad
\textbf{Ответственный:} Сотрудник C (верификация)

\subsection{Гравитация (эскиз)}
\begin{hypothesis}
Гравитационный эффект связан с градиентом интенсивности \VR\ в \SE, создаваемого ансамблем удалённых \EVO;
зависимость типа \(1/r^2\) возникает как сферическое ослабление возмущений и интегральная статистика импульсного обмена на \GP.
\end{hypothesis}

\subsection{Квантование (эскиз)}
\begin{hypothesis}
Дискретность устойчивых режимов связана с целочисленным замыканием циркуляции/геометрии канала
(самосогласованность статистики соударений на \GP). Простейшая запись для тонкого тора: \(R=n\,r\), \(n\in\mathbb{Z}\).
\end{hypothesis}

\subsection{Список проверок (пополняется)}
\begin{itemize}
\item Какие предсказания отличают модель от континуальных вакуумных теорий?
\item Какие измеримые параметры соответствуют \KP\ и интенсивности \VR?
\item Предельные режимы: когда формально воспроизводятся известные законы?
\end{itemize}

% sec_tests_C.tex
% Раздел C (Верификация и тесты): норматив, шлюзы, режимы допуска и витрина протокола
% Актуализировано под рабочий протокол C с QC1 и QC_SCREW_LOCK
% pdfLaTeX-safe: используем LaTeX-математику, без Unicode-математики в обычном тексте
%
% Требует в main.tex:
% \usepackage{amsmath,amssymb}
% \usepackage{tikz}
% \usetikzlibrary{arrows.meta,positioning,shapes.geometric}
% \usepackage{fancyvrb}
% \usepackage{hyperref}

\section{Верификация и тесты (C): норматив, шлюзы и условия допуска к интерпретации}
\label{sec:sec_tests_C}

\subsection{Почему потребовалось ужесточение протокола}

Верификационный контур C был ужесточён по двум независимым причинам.

Во-первых, локально-квазицилиндрическая запись удобна как технический инструмент расчёта, но сама по себе опасна: если использовать её без обязательной парной проверки с глобальной тороидальной постановкой, то можно перепутать \emph{координатный артефакт} с физическим эффектом. Поэтому локальный патч не может служить основанием для вывода без отдельной проверки геометрической инвариантности.

Во-вторых, дискретные шнековые режимы больше нельзя интерпретировать в старом непрерывном $q$-языке. В рабочем интерфейсе B$\to$C физический приоритет перенесён с непрерывного параметра $q$ на дискретную пару
\[
(m_A,m_S)
\]
и метку режима \texttt{mode\_S}. Параметр \texttt{q\_eff} оставлен только как вспомогательный локальный параметр наклона патча и не может служить самостоятельным основанием для физического вывода.

Именно поэтому в протокол C введены два обязательных шлюза:
\begin{itemize}
\item \textbf{QC1} --- защита от координатного артефакта в режиме \texttt{quasi\_cyl};
\item \textbf{QC\_SCREW\_LOCK} --- защита от произвольной непрерывной геометризации шнека.
\end{itemize}

Логика секции C строится в следующем порядке:
\begin{enumerate}
\item сначала фиксируются канонические запреты;
\item затем объясняется, почему \texttt{quasi\_cyl} опасен без пары \texttt{torus\_ref/quasi\_cyl\_ref};
\item затем объясняется, почему шнек нельзя интерпретировать через один \texttt{q\_eff};
\item далее вводятся QC1 и QC\_SCREW\_LOCK как охрана физического смысла;
\item только после этого перечисляются режимы, допускаемые к интерпретации.
\end{enumerate}

\subsection{Канонические запреты и язык наблюдения}

\textbf{(CANON \S16, \S17, норматив)} Прокси-набор $(K,G,A_{ij})$ используется только как
\emph{измерительный язык} состояния хаотической фазы у $\Sigma$ и \emph{не вводит новых сущностей}.
В тексте допустимы формулировки \emph{«наблюдаем прокси»} и запрещены формулировки вида
\emph{«действует поле $G/A$»}, \emph{«тензор $A$ толкает»}, \emph{«$A$ создаёт силу»}.

Рекомендуемая нормировка:
\[
K=n\langle v^2\rangle,\qquad
G=n\langle v\rangle,\qquad
P_{ij}=n\langle v_i v_j\rangle,\qquad
A_{ij}=P_{ij}-\frac13\mathrm{tr}(P)\delta_{ij}.
\]
На поверхности $\Sigma$ используем нормаль $n(x)$ и проектор на касательную плоскость
\[
\Pi_\top(\xi)=\xi-(\xi\cdot n)n.
\]

КУБК относится к балансу \emph{числа контактов} на $\Sigma$ и не является источником дрейфа.
Если приращения импульса изотропны, то устойчивый дрейф запрещён:
\[
A\equiv 0,\qquad \Pi_\top(G)\equiv 0 \quad \Longrightarrow \quad V=0.
\]
Здесь $V=0$ означает \textbf{векторный ноль}, а не только отсутствие продольной компоненты.

\subsection{Охранный тест T2 и нулевой коридор}

\textbf{(CANON \S4.3, \S7, \S17, норматив)} Тест T2 запрещает получать дрейф
из одной лишь $\lambda$-асимметрии. Если заявлено ``нет внешнего нарушителя'' и
входные прокси-метрики находятся в нулевом коридоре, то дрейф обязан оставаться нулевым.

Вводим метрики:
\[
v_{\mathrm{rms}}=\sqrt{\langle K/n\rangle_\Sigma},\qquad \bar n=\langle n\rangle_\Sigma,
\]
\[
\varepsilon_G=\frac{\|\Pi_\top(G)\|_\Sigma}{\bar n\,v_{\mathrm{rms}}},\qquad
\varepsilon_A=\frac{\|A\|_\Sigma}{\bar n\,v_{\mathrm{rms}}^2},
\]
\[
\varepsilon_V=\frac{\|V\|}{v_{\mathrm{rms}}},\qquad
\varepsilon_{V\parallel}=\frac{|V_\parallel|}{v_{\mathrm{rms}}}.
\]

\textbf{Авто-шлюз T2:}
если заявлено отсутствие внешнего нарушителя и
\[
(\varepsilon_G,\varepsilon_A)\approx 0,
\]
то обязано быть
\[
\varepsilon_V\approx 0.
\]

Красный флаг возникает, если $V\neq 0$ при $A\equiv 0$ и $\Pi_\top(G)\equiv 0$:
в этом случае дрейф фактически ``получен из $\lambda$'' или допущена ошибка постановки.

\subsection{Норматив тестов шнека T5--T6}

\textbf{(CANON \S17, норматив)} Фиксируется ось симметрии объекта $e_\parallel$,
измеряется продольная компонента дрейфа:
\[
V_\parallel = V\cdot e_\parallel .
\]

Обязательные проверки:
\[
V_\parallel(\chi)=-V_\parallel(-\chi),
\qquad
V_\parallel(s_{sh}=+1)=-V_\parallel(s_{sh}=-1)
\]
при прочих равных.

Для стоячего шнека при отсутствии внешнего нарушителя:
\[
\nabla K=0,\quad G=0,\quad A=0,\quad s_{sh}=0
\quad \Longrightarrow \quad
V_\parallel\approx 0.
\]

Метрики парных прогонов:
\[
\Delta_\chi=\frac{|V_\parallel(+1)+V_\parallel(-1)|}{v_{\mathrm{rms}}},
\qquad
\Delta_s=\frac{|V_\parallel(s_{sh}=+1)+V_\parallel(s_{sh}=-1)|}{v_{\mathrm{rms}}}.
\]

\subsection{QC1: защита от координатного артефакта \texttt{quasi\_cyl}}

Локальная запись \texttt{quasi\_cyl} допустима только как \emph{временный локальный патч}
той же канонической геометрии. Она не является новой физической сущностью и не заменяет
глобальную тороидальную/двуториальную постановку.

Поэтому для каждой канонической ситуации B обязан формировать минимум одну парную связку:
\begin{itemize}
\item \texttt{geometry\_mode=global\_torus}, \texttt{pair\_role=torus\_ref};
\item \texttt{geometry\_mode=quasi\_cyl}, \texttt{pair\_role=quasi\_cyl\_ref}.
\end{itemize}

При этом обязаны совпадать:
\begin{itemize}
\item \texttt{scenario},
\item \texttt{geometry\_id},
\item \texttt{pair\_id\_geom},
\item физические входы,
\item \texttt{chi},
\item \texttt{s\_sh},
\item режим внешнего нарушителя,
\item настройки допусков,
\end{itemize}
и различаться только:
\begin{itemize}
\item \texttt{geometry\_mode},
\item \texttt{pair\_role}.
\end{itemize}

\textbf{QC1 PASS:}
\begin{itemize}
\item охранные запреты T2, N1, N2, N3 не ломаются при переходе к \texttt{quasi\_cyl};
\item новый дрейф не возникает только из выбора координатного патча;
\item допустимость/недопустимость D1 совпадает в паре \texttt{torus\_ref/quasi\_cyl\_ref}.
\end{itemize}

\textbf{QC1 FAIL:}
\begin{itemize}
\item запреты нарушаются только в \texttt{quasi\_cyl};
\item новый эффект есть только в \texttt{quasi\_cyl}, а в \texttt{global\_torus} исчезает.
\end{itemize}

В этом случае серия маркируется как
\[
\texttt{coordinate\_artifact\_suspected},
\]
а физическая интерпретация откладывается.

\subsection{QC\_SCREW\_LOCK: защита от произвольной непрерывной геометризации шнека}

В новой рабочей схеме B$\to$C физика шнека читается через дискретную пару
\[
(m_A,m_S),
\]
где:
\begin{itemize}
\item $m_A\in \mathbb{Z}_{>0}$ --- квантовое число канальной навивки;
\item $m_S\in\{1,2,4\}$ --- базовая шнековая мода;
\item \texttt{mode\_S} --- метка режима:
\[
\texttt{mode\_S}\in\{\texttt{S1},\texttt{S2},\texttt{S4},\texttt{S*}\}.
\]
\end{itemize}

Базовые режимы:
\[
\texttt{S1},\ \texttt{S2},\ \texttt{S4}.
\]
Режим \texttt{S*} допускается только как \emph{производный} режим с явной маркировкой производности.

Дополнительно B обязан передавать:
\begin{itemize}
\item \texttt{hA\_definition},
\item \texttt{hS\_definition},
\item \texttt{A\_S\_orth\_candidate},
\item \texttt{q\_eff},
\item \texttt{mode\_lock\_pass}.
\end{itemize}

\textbf{Критический смысл:} \texttt{q\_eff} не является главным физическим носителем шнека.
Он допустим только как вспомогательный локальный параметр патча.

\textbf{QC\_SCREW\_LOCK PASS:}
\begin{itemize}
\item серия читается через согласованную пару $(m_A,m_S)$;
\item базовые режимы принадлежат множеству \texttt{S1/S2/S4};
\item \texttt{S*} указан как производный режим;
\item \texttt{h\_A} и \texttt{h\_S} разведены явно;
\item \texttt{q\_eff} не используется как физическая основа режима;
\item \texttt{mode\_lock\_pass=true}.
\end{itemize}

\textbf{QC\_SCREW\_LOCK FAIL:}
\begin{itemize}
\item шнековая мода подана как произвольная непрерывная деформация;
\item связь с $m_A$ отсутствует;
\item режим вне базы \texttt{S1/S2/S4} заявлен без указания производности;
\item физическая интерпретация даётся при \texttt{mode\_lock\_pass=false};
\item \texttt{q\_eff} используется как главный аргумент физического вывода.
\end{itemize}

\subsection{Базовые моды S1/S2/S4 и условия допуска к интерпретации}

Для каждого нового шнекового сценария C требует:
\begin{enumerate}
\item парность \texttt{global\_torus} $\leftrightarrow$ \texttt{quasi\_cyl};
\item сравнение базовых мод \texttt{S1}, \texttt{S2}, \texttt{S4};
\item сравнение ``базовая мода $\leftrightarrow$ производная мода'';
\item сравнение ``согласованная пара $(m_A,m_S)$ $\leftrightarrow$ искусственно рассогласованная пара''.
\end{enumerate}

Физическая интерпретация шнековой серии допустима только если одновременно:
\begin{itemize}
\item \texttt{local\_patch\_valid=true};
\item \texttt{patch\_valid=true} (если используется патч-язык);
\item \texttt{QC1\_pass=true};
\item \texttt{QC\_SCREW\_LOCK\_pass=true};
\item \texttt{mode\_lock\_pass=true};
\item пройден T2;
\item выполнены T5--T6;
\item пройдены N1--N3;
\item D1 заявляется только в допустимом режиме.
\end{itemize}

Если хотя бы одно из условий нарушено, C обязан выставить:
\[
\texttt{FAIL\_INTERPRETATION}
\]
и/или соответствующий красный флаг.

\subsection{Охранники N1--N3 и режим D1 в новой схеме}

Перед тем, как считать режим D1 физически интерпретируемым, обязательны охранники:

\begin{itemize}
\item \textbf{N1 (базовая тороидальная ссылка):}
нет допустимого дискретно-согласованного пространственного режима; локальный tor-вклад и вращение допустимы,
но интегрально обязано быть
\[
V\approx 0.
\]

\item \textbf{N2 (без tor-моды):}
\[
s_{sh}=0 \quad \Longrightarrow \quad
j_{\mathrm{tor}}^{\mathrm{loc}}=0
\Longrightarrow
j_{\mathrm{sp}}^{\mathrm{loc}}=0
\Longrightarrow
V=0.
\]

\item \textbf{N3 (без геометрического рычага):}
\[
w_{\mathrm{asym}}\approx 0 \quad \Longrightarrow \quad V\approx 0.
\]
\end{itemize}

\textbf{D1 (дискретно-согласованный двуторий со шнеком)} допускается только если:
\begin{itemize}
\item предъявлена допустимая пара $(m_A,m_S)$;
\item \texttt{mode\_S}$\in\{\texttt{S1},\texttt{S2},\texttt{S4}\}$ или явно обоснованная производная \texttt{S*};
\item $w_{\mathrm{asym}}>0$;
\item $s_{sh}\neq 0$.
\end{itemize}

Только здесь допускается
\[
V_\parallel\neq 0,
\]
но с обязательным выполнением норм CANON \S17:
\[
V_\parallel(\chi)=-V_\parallel(-\chi),
\qquad
V_\parallel(s_{sh}=+1)=-V_\parallel(s_{sh}=-1).
\]

\subsection{Почему в витрине должны быть две поточные схемы}

В этой секции намеренно сохраняются \textbf{две} визуальные схемы.

Первая схема --- \emph{историческая}. Она нужна не как действующий норматив, а как
документированная фиксация прошлого этапа, в котором C строил объяснение через
непрерывную геометризацию двутория при помощи $q$, $h(x)$ и $w(x)$.
Эта схема показывает, откуда произошёл ранний язык охранников N1--N3 и D1.

Вторая схема --- \emph{актуальная}. Она нужна как действующий поток проверки в текущем
рабочем протоколе C, где центральны уже не параметры $q/h(x)$, а:
\begin{itemize}
\item \texttt{geometry\_mode},
\item временный локальный патч,
\item QC1,
\item дискретная пара $(m_A,m_S)$,
\item \texttt{mode\_S},
\item QC\_SCREW\_LOCK,
\item и только затем --- физическая интерпретация.
\end{itemize}

Сохранение обеих схем нужно для честной демонстрации всех шагов эволюции протокола:
мы не скрываем прежний каркас, но и не выдаём его за действующий.

\subsection{Историческая поточная схема C (ранняя q-центричная витрина)}

\begin{figure}[p]
\centering
\resizebox{\textwidth}{0.90\textheight}{%
\begin{tikzpicture}[
  font=\normalsize,
  node distance=9mm and 24mm,
  box/.style={rectangle, rounded corners=3pt, draw=black!55, thick, align=center,
              minimum width=0.78\linewidth, inner sep=8pt},
  sbox/.style={rectangle, rounded corners=3pt, draw=black!55, thick, align=center,
               minimum width=0.36\linewidth, inner sep=8pt},
  decision/.style={diamond, aspect=2.25, draw=black!55, thick, align=center,
                   inner sep=3pt, text width=0.44\linewidth},
  arr/.style={-Latex, thick, draw=black!70}
]

\node[box, fill=blue!10] (geom) {%
$\Sigma$: граничная поверхность, параметризация $X(u,v)$\\
$n(x)$: нормаль на $\Sigma$,\quad $e_{\parallel}$: ось (CANON \S17),\quad $V_{\parallel}=V\cdot e_{\parallel}$\\
паспорт: $\chi\in\{-1,+1\}$,\quad $s_{sh}\in\{-1,0,+1\}$ (бег $\rightarrow$/стоячий/бег $\leftarrow$)
};

\node[box, fill=cyan!10, below=of geom] (geo) {%
\textbf{Геометризация двутория (B$\to$C, КАНДИДАТ):}\\
$q\in\mathbb{R}\setminus\{0\}$ (q=0: тонкий тор; q$\neq$0: двуторий),\ \ $w(x)$: вес внешн./внутр. рычага\\
\texttt{w\_mode $\in$ \{binary, hybrid, ...\}},\quad \texttt{w\_asym} — метрика несимметрии $w(x)$\\
Норма каркаса: дрейф двутория запрещён без (q$\neq$0 и w\_asym$>$0) и без tor-моды (s\_sh$\neq$0)
};

\node[box, fill=orange!12, below=of geo] (proxies) {%
Прокси на $\Sigma$ (CANON \S16):\\
$K=n\langle v^2\rangle,\quad G=n\langle v\rangle,\quad P_{ij}=n\langle v_i v_j\rangle$\\
$A_{ij}=P_{ij}-\frac13\mathrm{tr}(P)\delta_{ij},\quad \Pi_{\top}(G)=G-(G\!\cdot\!n)n$\\
\textbf{Прокси — ярлыки измерения, не сущности и не «поля».}
};

\node[decision, fill=white, below=of proxies] (dec) {Расчёт $j$ на $\Sigma$};

\node[sbox, fill=violet!12, below left=of dec, yshift=+4mm] (ext) {%
\textbf{Внешний}\\[2pt]
$j^{ext}(x)=F_{ext}(K,G,A;\Sigma)$\\
$\nabla K$ / приходящая ВР\\
(перекосы статистики на $\Sigma$)
};

\node[sbox, fill=violet!12, below right=of dec, yshift=+4mm] (intr) {%
\textbf{Внутренний (tor$\to$sp)}\\[2pt]
$g_{\top}=\Pi_{\top}(G),\quad a_{\top}=\Pi_{\top}(A n)$\\
$j_{tor}=\chi\,s_{sh}(\alpha_1\,n\times g_{\top}+\alpha_2\,n\times a_{\top})$\\
$j_{sp}=\beta(K)\,w(x)\,(j_{tor}\!\cdot\!t_\theta)\,h(x)$\\
$j_{int}=j_{tor}+j_{sp}$\\
\textbf{s\_sh=0 $\Rightarrow$ $j_{tor}=j_{sp}=0$}
};

\node[box, fill=green!12, below=30mm of dec] (sum) {%
$j(x)=j^{ext}(x)+j^{int}(x)$\\
локальная импульсная анизотропия контактов на $\Sigma$
};

\node[box, fill=purple!10, below=of sum] (resp) {%
Интегральный отклик:\\
$V=\mu_T\int_{\Sigma} j\,dS,\quad V_{\parallel}=V\cdot e_{\parallel}$\\
$\Omega=\tilde{\mu}_R\int_{\Sigma}(r-r_0)\times j\,dS$\\
$J_{int,\parallel}=\Big(\int_{\Sigma} j^{int}\,dS\Big)\cdot e_{\parallel}$
};

\node[box, fill=yellow!18, below=of resp] (ver) {%
Верификация / нормы:\\
$v_{rms}=\sqrt{\langle K/n\rangle_{\Sigma}},\quad \bar n=\langle n\rangle_{\Sigma}$\\
$\varepsilon_G=\|\Pi_{\top}(G)\|_{\Sigma}/(\bar n\,v_{rms}),\quad
\varepsilon_A=\|A\|_{\Sigma}/(\bar n\,v_{rms}^2)$\\
$\varepsilon_V=\|V\|/v_{rms},\quad \varepsilon_{V\parallel}=|V_{\parallel}|/v_{rms}$\\
\textbf{T2:} $(\varepsilon_G,\varepsilon_A)\approx 0 \Rightarrow \varepsilon_V\approx 0$\\
\textbf{T5--T6:} $\chi$-нечётность $V_{\parallel}$ и реверс $s_{sh}$; $s_{sh}=0 \Rightarrow V_{\parallel}\approx 0$\\[4pt]
\textbf{Охранники геометризации (перед D1 обязательно):}\\
N1: $q=0,\ s_{sh}\neq 0 \Rightarrow \Omega\neq 0$ допустимо, но $V\approx 0$\\
N2: $q\neq 0,\ s_{sh}=0 \Rightarrow V=0$\\
N3: $q\neq 0,\ w\_{asym}\approx 0 \Rightarrow V\approx 0$\\
\textbf{D1:} $q\neq 0,\ w\_{asym}>0,\ s_{sh}\neq 0 \Rightarrow V_{\parallel}\neq 0$ + нормы CANON \S17
};

\draw[arr] (geom) -- (geo);
\draw[arr] (geo) -- (proxies);
\draw[arr] (proxies) -- (dec);
\draw[arr] (dec) -- (ext);
\draw[arr] (dec) -- (intr);
\draw[arr] (ext.south)  -- ++(0,-10mm) |- (sum.west);
\draw[arr] (intr.south) -- ++(0,-10mm) |- (sum.east);
\draw[arr] (sum) -- (resp);
\draw[arr] (resp) -- (ver);

\end{tikzpicture}%
}
\caption{Историческая поточная схема (C): ранняя витрина через непрерывную геометризацию двутория $(q,w,h)$ и переход $tor\to sp$. Схема сохраняется как документированная версия прежнего этапа развития протокола.}
\label{fig:C_flow_tikz_old}
\end{figure}

\subsection{Актуальная поточная схема C (действующий протокол)}

Новая схема вводится не для косметической замены обозначений, а для изменения самой логики проверки.
Если историческая схема показывала, как C раньше описывал внутренний механизм через непрерывную геометризацию,
то актуальная схема показывает, как C теперь проверяет \emph{допустимость физического вывода}.

В новой логике:
\begin{enumerate}
\item локальный патч признаётся только временным инструментом расчёта;
\item затем обязательна геометрическая проверка QC1;
\item затем обязательна модовая проверка QC\_SCREW\_LOCK;
\item только после этого допускается интерпретация результата как физического.
\end{enumerate}

Именно поэтому в новой схеме:
\begin{itemize}
\item появился шаг \emph{временного локального патча};
\item появился явный шлюз \textbf{QC1};
\item появился явный шлюз \textbf{QC\_SCREW\_LOCK};
\item параметр $q$ ушёл с роли главного физического носителя шнека;
\item интерпретация вынесена в последний блок потока.
\end{itemize}

\begin{figure}[p]
\centering
% Новый блок на ТОЧНО ТОМ ЖЕ геометрическом каркасе, что и old:
% меняется только содержание узлов, а не внешняя компоновка.
\resizebox{\textwidth}{0.90\textheight}{%
\begin{tikzpicture}[
  font=\normalsize,
  node distance=9mm and 24mm,
  box/.style={rectangle, rounded corners=3pt, draw=black!55, thick, align=center,
              minimum width=0.78\linewidth, inner sep=8pt},
  sbox/.style={rectangle, rounded corners=3pt, draw=black!55, thick, align=center,
               minimum width=0.36\linewidth, inner sep=8pt},
  decision/.style={diamond, aspect=2.25, draw=black!55, thick, align=center,
                   inner sep=3pt, text width=0.44\linewidth},
  arr/.style={-Latex, thick, draw=black!70}
]

% --- 1) Геометрия
\node[box, fill=blue!10] (geom) {%
$\Sigma$: граничная поверхность, параметризация $X(u,v)$\\
$n(x)$: нормаль на $\Sigma$,\quad $e_{\parallel}$: ось (CANON \S17),\quad $V_{\parallel}=V\cdot e_{\parallel}$\\
паспорт: $\chi\in\{-1,+1\}$,\quad $s_{sh}\in\{-1,0,+1\}$,\quad \texttt{geometry\_mode $\in$ \{global\_torus, quasi\_cyl\}}
};

% --- 2) Временный локальный патч
\node[box, fill=cyan!10, below=of geom] (geo) {%
\textbf{Временный локальный патч и геометрический режим}\\
\texttt{global\_torus}: глобальная тороидальная / двуториальная постановка\\
\texttt{quasi\_cyl}: только локальный патч записи $\Sigma_{loc}\subset\Sigma$\\
\texttt{local\_patch\_valid}, \texttt{patch\_valid}, \texttt{curvature\_scale\_ratio}, \texttt{curvature\_ratio}\\
\textbf{Норма:} \texttt{quasi\_cyl} не является новой физической сущностью
};

% --- 3) Прокси
\node[box, fill=orange!12, below=of geo] (proxies) {%
Прокси на $\Sigma$ (CANON \S16):\\
$K=n\langle v^2\rangle,\quad G=n\langle v\rangle,\quad P_{ij}=n\langle v_i v_j\rangle$\\
$A_{ij}=P_{ij}-\frac13\mathrm{tr}(P)\delta_{ij},\quad \Pi_{\top}(G)=G-(G\!\cdot\!n)n$\\
\textbf{Прокси --- ярлыки измерения, не сущности и не «поля».}
};

% --- 4) Ромб: QC1
\node[decision, fill=white, below=of proxies] (dec) {%
\textbf{QC1}\\
геометрическая инвариантность\\
\texttt{torus\_ref} $\leftrightarrow$ \texttt{quasi\_cyl\_ref}
};

% --- 5) Внешний
\node[sbox, fill=violet!12, below left=of dec, yshift=+4mm] (ext) {%
\textbf{Внешний}\\[2pt]
$j^{ext}(x)=F_{ext}(K,G,A;\Sigma)$\\
$\nabla K$ / приходящая ВР\\
внешний перекос статистики на $\Sigma$
};

% --- 6) Внутренний: новая модовая схема
\node[sbox, fill=violet!12, below right=of dec, yshift=+4mm] (intr) {%
\textbf{Внутренний (дискретно-модовый)}\\[2pt]
$m_A\in\mathbb{Z}_{>0},\quad m_S\in\{1,2,4\}$,\quad \texttt{mode\_S $\in$ \{S1,S2,S4,S*\}}\\
$h_A=h_A(m_A;x),\quad h_S=h_S(m_S,\chi,s_{sh};x)$\\
$j_{tor}^{loc}=\chi\,s_{sh}\,\mathcal{T}(m_A;K,g_{\top},a_{\top},n)$\\
$j_{sp}^{loc}=\beta(K;m_A,m_S)\,w_{loc}(x)\,\Gamma_{A,S}(x)\,h_S(x)$\\
\textbf{\texttt{q\_eff} --- только вспомогательный параметр патча}
};

% --- 7) Сумма / модовый шлюз
\node[box, fill=green!12, below=30mm of dec] (sum) {%
\textbf{QC\_SCREW\_LOCK и суммарный отклик}\\
база: \texttt{S1/S2/S4},\quad \texttt{S*} только как производная,\quad \texttt{mode\_lock\_pass=true}\\
$j(x)=j^{ext}(x)+j^{int}(x)$\\
локальная импульсная анизотропия контактов на $\Sigma$
};

% --- 8) Интегральный отклик
\node[box, fill=purple!10, below=of sum] (resp) {%
Интегральный отклик:\\
$V=\mu_T\int_{\Sigma} j\,dS,\quad V_{\parallel}=V\cdot e_{\parallel}$\\
$\Omega=\tilde{\mu}_R\int_{\Sigma}(r-r_0)\times j\,dS$\\
$J_{int,\parallel}=\Big(\int_{\Sigma} j^{int}\,dS\Big)\cdot e_{\parallel}$
};

% --- 9) Нормы/шлюзы + N1–N3 + D1
\node[box, fill=yellow!18, below=of resp] (ver) {%
Верификация / нормы:\\
$v_{rms}=\sqrt{\langle K/n\rangle_{\Sigma}},\quad \bar n=\langle n\rangle_{\Sigma}$\\
$\varepsilon_G=\|\Pi_{\top}(G)\|_{\Sigma}/(\bar n\,v_{rms}),\quad
\varepsilon_A=\|A\|_{\Sigma}/(\bar n\,v_{rms}^2)$\\
$\varepsilon_V=\|V\|/v_{rms},\quad \varepsilon_{V\parallel}=|V_{\parallel}|/v_{rms}$\\
\textbf{T2:} $(\varepsilon_G,\varepsilon_A)\approx 0 \Rightarrow \varepsilon_V\approx 0$\\
\textbf{T5--T6:} $\chi$-нечётность $V_{\parallel}$ и реверс $s_{sh}$; $s_{sh}=0 \Rightarrow V_{\parallel}\approx 0$\\[4pt]
\textbf{Охранники перед D1:}\\
N1: базовая тороидальная ссылка $\Rightarrow V\approx 0$\\
N2: $s_{sh}=0 \Rightarrow j_{tor}^{loc}=0 \Rightarrow j_{sp}^{loc}=0 \Rightarrow V=0$\\
N3: $w\_{asym}\approx 0 \Rightarrow V\approx 0$\\
\textbf{D1:} допустим только при корректной паре $(m_A,m_S)$ и режиме \texttt{S1/S2/S4} или производной \texttt{S*}
};

% --- Стрелки
\draw[arr] (geom) -- (geo);
\draw[arr] (geo) -- (proxies);
\draw[arr] (proxies) -- (dec);
\draw[arr] (dec) -- (ext);
\draw[arr] (dec) -- (intr);

% Входы в sum по серединам граней (как в old):
\draw[arr] (ext.south)  -- ++(0,-10mm) |- (sum.west);
\draw[arr] (intr.south) -- ++(0,-10mm) |- (sum.east);

\draw[arr] (sum) -- (resp);
\draw[arr] (resp) -- (ver);

\end{tikzpicture}%
}
\caption{Актуальная поточная схема (C): сохранён прежний внешний каркас витрины, но внутренняя логика обновлена под режимы \texttt{global\_torus/quasi\_cyl}, геометрический шлюз QC1, дискретно-модовую схему шнека и модовый шлюз QC\_SCREW\_LOCK.}
\label{fig:C_flow_tikz_new}
\end{figure}

\FloatBarrier

\subsection{Реестр красных флагов}

\begin{itemize}
\item дрейф появляется при изотропном фоне без нарушителя;
\item дрейф ``объясняется'' асимметрией числа контактов, а не анизотропией $\Delta p$;
\item $\Phi$ интерпретируется как поток СЭ;
\item $(K,G,A_{ij})$ трактуются как действующие поля/сущности;
\item \texttt{quasi\_cyl} трактуется как отдельная физическая сущность или глобальная форма ЭВО;
\item новый эффект существует только в \texttt{quasi\_cyl};
\item \texttt{q\_eff} используется как главный физический носитель шнека;
\item режим вне базы \texttt{S1/S2/S4} интерпретируется без явной производности;
\item физическая интерпретация дана при \texttt{mode\_lock\_pass=false};
\item D1 заявлен без прохождения N1--N3, QC1 и QC\_SCREW\_LOCK.
\end{itemize}


% =========================================================
\appendix
\section{Журнал изменений}
\subsection*{2026-02-11}
\begin{itemize}
\item Создан базовый LaTeX-каркас.
\item В канон внесено условие целостности хаотической фазы: \(l\le d_A\), \(\mean{l}=0.5\,d_A\).
\end{itemize}

\subsection*{2026-02-17}
\begin{itemize}
\item Канон обновлён: добавлены КУБК, тонкий тор, динамика шнека, симметрия дрейфа, двуторий как режим ``вращение$\to$дрейф''.
\item Математическое ядро вынесено в отдельный файл и подключено через \texttt{\textbackslash input}.
\end{itemize}

\subsection*{2026-02-19}
\begin{itemize}
\item Канон расширён до \textbf{1--17}: добавлены \S16 (прокси \((K,G,A_{ij})\) как измерительные ярлыки \(\KP/\VR\), не новые сущности)
и \S17 (норматив тестов шнека: \(\chi\)-нечётность \(V_{\parallel}\), реверс бега шнека).
\item В \S6 устранена терминологическая ловушка ``потоки к поверхности'': \(\Phi\) закреплена как \textbf{интенсивность контактов} на \(\Sigma\).
\end{itemize}

\subsection*{2026-03-09}

\begin{itemize}
\item Уточнена логика текущего шага проекта: квазицилиндрический патч $\Sigma_{\mathrm{loc}}$
введён не как замена глобальной тороидальной/двуториальной геометрии ЭВО,
а как временное локальное сужение задачи, выбранное для более явного проявления
механизмов дрейфа и ротации через контактную кинетику на граничной поверхности $\Sigma$.

\item Математический каркас B переведён с непрерывной q-центричной геометризации
на дискретно-согласованную схему шнека:
$m_A \in \mathbb{Z}_{>0}$ --- квантовое число канальной навивки,
$m_S \in \{1,2,4\}$ --- базовая пространственная мода шнека,
$S1/S2/S4$ --- базовые режимы,
$S^\*$ --- только производные режимы.

\item В локальном патч-языке разведены два касательных направления:
$h_A$ --- директор траектории Амеров,
$h_S$ --- директор шнековой моды ГП.
Тем самым снята прежняя перегрузка непрерывного параметра $q$;
он сохранён только как вспомогательный эффективный локальный параметр $q_{\mathrm{eff}}$.

\item Верификационный протокол C ужесточён:
для всех критичных серий обязательны
парность \texttt{global\_torus $\leftrightarrow$ quasi\_cyl},
геометрический шлюз \texttt{QC1}
и модовый шлюз \texttt{QC\_SCREW\_LOCK}.
Без их прохождения физическая интерпретация запрещена.

\item Витрина теперь должна отражать последовательность шагов проекта:
глобальная цель $\rightarrow$ тонкий тор $\rightarrow$ локальный патч $\Sigma_{\mathrm{loc}}$
$\rightarrow$ разведение $h_A/h_S$
$\rightarrow$ дискретные моды $S1/S2/S4$
$\rightarrow$ возврат к глобальному интегралу по $\Sigma$.
\end{itemize}

\end{document}
